\documentclass[11pt,a4paper]{article}
% ============================================
% GIFT v3.4 — Shared Preamble
% ============================================
% Usage: % ============================================
% GIFT v3.4 — Shared Preamble
% ============================================
% Usage: % ============================================
% GIFT v3.4 — Shared Preamble
% ============================================
% Usage: \input{gift_preamble} after \documentclass

% ENCODING & FONTS
\usepackage[utf8]{inputenc}
\usepackage[T1]{fontenc}
\usepackage{lmodern}

% PAGE LAYOUT (Golden Ratio)
\usepackage[margin=1.618cm, top=2.618cm, bottom=2.618cm]{geometry}

% ESSENTIAL PACKAGES
\usepackage{float}
\usepackage{caption}
\usepackage{subcaption}
\usepackage{setspace}
\usepackage{fancyhdr}
\usepackage{xcolor}
\usepackage{hyperref}
\usepackage{csquotes}
\usepackage{amsmath}
\usepackage{amssymb}
\usepackage{amsthm}
\usepackage{mathtools}
\usepackage{booktabs}
\usepackage{longtable}
\usepackage{array}
\usepackage{multirow}
\usepackage{tikz}
\usepackage{graphicx}
\usepackage{listings}
\usepackage{enumitem}
% Unicode fallbacks (pandoc artifacts that survive markdown→LaTeX)
\DeclareUnicodeCharacter{00B0}{\ensuremath{^\circ}}
\DeclareUnicodeCharacter{221A}{\ensuremath{\sqrt{\,}}}
\DeclareUnicodeCharacter{2713}{\ensuremath{\checkmark}}
\DeclareUnicodeCharacter{2717}{\ensuremath{\times}}
\DeclareUnicodeCharacter{2070}{\ensuremath{^{0}}}
\DeclareUnicodeCharacter{00B9}{\ensuremath{^{1}}}
\DeclareUnicodeCharacter{00B2}{\ensuremath{^{2}}}
\DeclareUnicodeCharacter{00B3}{\ensuremath{^{3}}}
\DeclareUnicodeCharacter{2074}{\ensuremath{^{4}}}
\DeclareUnicodeCharacter{2075}{\ensuremath{^{5}}}
\DeclareUnicodeCharacter{2076}{\ensuremath{^{6}}}
\DeclareUnicodeCharacter{2077}{\ensuremath{^{7}}}
\DeclareUnicodeCharacter{2078}{\ensuremath{^{8}}}
\DeclareUnicodeCharacter{2079}{\ensuremath{^{9}}}
\DeclareUnicodeCharacter{2071}{\ensuremath{^{i}}}
\DeclareUnicodeCharacter{2080}{\ensuremath{_{0}}}
\DeclareUnicodeCharacter{2081}{\ensuremath{_{1}}}
\DeclareUnicodeCharacter{2082}{\ensuremath{_{2}}}
\DeclareUnicodeCharacter{2083}{\ensuremath{_{3}}}
\DeclareUnicodeCharacter{2084}{\ensuremath{_{4}}}
\DeclareUnicodeCharacter{2085}{\ensuremath{_{5}}}
\DeclareUnicodeCharacter{2086}{\ensuremath{_{6}}}
\DeclareUnicodeCharacter{2087}{\ensuremath{_{7}}}
\DeclareUnicodeCharacter{2088}{\ensuremath{_{8}}}
\DeclareUnicodeCharacter{2089}{\ensuremath{_{9}}}
\DeclareUnicodeCharacter{208A}{\ensuremath{_{+}}}
\DeclareUnicodeCharacter{208B}{\ensuremath{_{-}}}
\DeclareUnicodeCharacter{2192}{\ensuremath{\to}}
\DeclareUnicodeCharacter{2502}{|}
\DeclareUnicodeCharacter{25BC}{\ensuremath{\blacktriangledown}}
\DeclareUnicodeCharacter{2500}{\ensuremath{-}}
\DeclareUnicodeCharacter{2514}{\ensuremath{\llcorner}}
\DeclareUnicodeCharacter{251C}{\ensuremath{\vdash}}
\DeclareUnicodeCharacter{2026}{\ldots}
\DeclareUnicodeCharacter{2208}{\ensuremath{\in}}
\DeclareUnicodeCharacter{2229}{\ensuremath{\cap}}
\DeclareUnicodeCharacter{222A}{\ensuremath{\cup}}
\DeclareUnicodeCharacter{2295}{\ensuremath{\oplus}}
\DeclareUnicodeCharacter{2297}{\ensuremath{\otimes}}
\DeclareUnicodeCharacter{2245}{\ensuremath{\cong}}
\DeclareUnicodeCharacter{2248}{\ensuremath{\approx}}
\DeclareUnicodeCharacter{2260}{\ensuremath{\neq}}
\DeclareUnicodeCharacter{2261}{\ensuremath{\equiv}}
\DeclareUnicodeCharacter{2264}{\ensuremath{\leq}}
\DeclareUnicodeCharacter{2265}{\ensuremath{\geq}}
\DeclareUnicodeCharacter{00D7}{\ensuremath{\times}}
\DeclareUnicodeCharacter{00B1}{\ensuremath{\pm}}
\DeclareUnicodeCharacter{2225}{\ensuremath{\|}}
% Greek letters used in body text (pandoc passes them through unchanged)
\DeclareUnicodeCharacter{03B1}{\ensuremath{\alpha}}
\DeclareUnicodeCharacter{03B2}{\ensuremath{\beta}}
\DeclareUnicodeCharacter{03B3}{\ensuremath{\gamma}}
\DeclareUnicodeCharacter{03B4}{\ensuremath{\delta}}
\DeclareUnicodeCharacter{03B5}{\ensuremath{\varepsilon}}
\DeclareUnicodeCharacter{03B6}{\ensuremath{\zeta}}
\DeclareUnicodeCharacter{03B7}{\ensuremath{\eta}}
\DeclareUnicodeCharacter{03B8}{\ensuremath{\theta}}
\DeclareUnicodeCharacter{03B9}{\ensuremath{\iota}}
\DeclareUnicodeCharacter{03BA}{\ensuremath{\kappa}}
\DeclareUnicodeCharacter{03BB}{\ensuremath{\lambda}}
\DeclareUnicodeCharacter{03BC}{\ensuremath{\mu}}
\DeclareUnicodeCharacter{03BD}{\ensuremath{\nu}}
\DeclareUnicodeCharacter{03BE}{\ensuremath{\xi}}
\DeclareUnicodeCharacter{03C0}{\ensuremath{\pi}}
\DeclareUnicodeCharacter{03C1}{\ensuremath{\rho}}
\DeclareUnicodeCharacter{03C3}{\ensuremath{\sigma}}
\DeclareUnicodeCharacter{03C4}{\ensuremath{\tau}}
\DeclareUnicodeCharacter{03C5}{\ensuremath{\upsilon}}
\DeclareUnicodeCharacter{03C6}{\ensuremath{\varphi}}
\DeclareUnicodeCharacter{03C7}{\ensuremath{\chi}}
\DeclareUnicodeCharacter{03C8}{\ensuremath{\psi}}
\DeclareUnicodeCharacter{03C9}{\ensuremath{\omega}}
\DeclareUnicodeCharacter{0394}{\ensuremath{\Delta}}
\DeclareUnicodeCharacter{0398}{\ensuremath{\Theta}}
\DeclareUnicodeCharacter{039B}{\ensuremath{\Lambda}}
\DeclareUnicodeCharacter{03A0}{\ensuremath{\Pi}}
\DeclareUnicodeCharacter{03A3}{\ensuremath{\Sigma}}
\DeclareUnicodeCharacter{03A6}{\ensuremath{\Phi}}
\DeclareUnicodeCharacter{03A8}{\ensuremath{\Psi}}
\DeclareUnicodeCharacter{03A9}{\ensuremath{\Omega}}
\DeclareUnicodeCharacter{2200}{\ensuremath{\forall}}
\DeclareUnicodeCharacter{2203}{\ensuremath{\exists}}
\DeclareUnicodeCharacter{2207}{\ensuremath{\nabla}}
\DeclareUnicodeCharacter{2202}{\ensuremath{\partial}}
\DeclareUnicodeCharacter{2218}{\ensuremath{\circ}}
\DeclareUnicodeCharacter{2032}{\ensuremath{\prime}}
\DeclareUnicodeCharacter{2113}{\ensuremath{\ell}}
\DeclareUnicodeCharacter{210F}{\ensuremath{\hbar}}
\DeclareUnicodeCharacter{27E8}{\ensuremath{\langle}}
\DeclareUnicodeCharacter{27E9}{\ensuremath{\rangle}}
\DeclareUnicodeCharacter{2308}{\ensuremath{\lceil}}
\DeclareUnicodeCharacter{2309}{\ensuremath{\rceil}}
\DeclareUnicodeCharacter{230A}{\ensuremath{\lfloor}}
\DeclareUnicodeCharacter{230B}{\ensuremath{\rfloor}}
\DeclareUnicodeCharacter{222B}{\ensuremath{\int}}
\DeclareUnicodeCharacter{2211}{\ensuremath{\sum}}
\DeclareUnicodeCharacter{220F}{\ensuremath{\prod}}
\DeclareUnicodeCharacter{221E}{\ensuremath{\infty}}
\DeclareUnicodeCharacter{00B5}{\ensuremath{\mu}}
\DeclareUnicodeCharacter{2236}{\ensuremath{\colon}}
\DeclareUnicodeCharacter{2192}{\ensuremath{\to}}
\DeclareUnicodeCharacter{21A6}{\ensuremath{\mapsto}}
\DeclareUnicodeCharacter{21D2}{\ensuremath{\Rightarrow}}
\DeclareUnicodeCharacter{21D4}{\ensuremath{\Leftrightarrow}}
\DeclareUnicodeCharacter{2282}{\ensuremath{\subset}}
\DeclareUnicodeCharacter{2286}{\ensuremath{\subseteq}}
\DeclareUnicodeCharacter{2287}{\ensuremath{\supseteq}}
\DeclareUnicodeCharacter{222B}{\ensuremath{\int}}
\DeclareUnicodeCharacter{2205}{\ensuremath{\emptyset}}
\DeclareUnicodeCharacter{2299}{\ensuremath{\odot}}
\DeclareUnicodeCharacter{2220}{\ensuremath{\angle}}
\DeclareUnicodeCharacter{2135}{\ensuremath{\aleph}}
\DeclareUnicodeCharacter{207B}{\ensuremath{^{-}}}
\DeclareUnicodeCharacter{207A}{\ensuremath{^{+}}}
\DeclareUnicodeCharacter{220E}{\ensuremath{\blacksquare}}
\DeclareUnicodeCharacter{2061}{}
\DeclareUnicodeCharacter{2062}{}
\DeclareUnicodeCharacter{2063}{}
\DeclareUnicodeCharacter{2009}{\,}
\DeclareUnicodeCharacter{200A}{\,}
\DeclareUnicodeCharacter{200B}{}
\DeclareUnicodeCharacter{2212}{\ensuremath{-}}
\DeclareUnicodeCharacter{2194}{\ensuremath{\leftrightarrow}}
\DeclareUnicodeCharacter{2099}{\ensuremath{_{n}}}
\DeclareUnicodeCharacter{2098}{\ensuremath{_{m}}}
\DeclareUnicodeCharacter{2096}{\ensuremath{_{k}}}
\DeclareUnicodeCharacter{2095}{\ensuremath{_{h}}}
\DeclareUnicodeCharacter{2090}{\ensuremath{_{a}}}
\DeclareUnicodeCharacter{2091}{\ensuremath{_{e}}}
\DeclareUnicodeCharacter{2093}{\ensuremath{_{x}}}
\DeclareUnicodeCharacter{1D62}{\ensuremath{_{i}}}
\DeclareUnicodeCharacter{2C7C}{\ensuremath{_{j}}}
\DeclareUnicodeCharacter{2C7B}{\ensuremath{^{e}}}
\DeclareUnicodeCharacter{1D52}{\ensuremath{^{o}}}
\DeclareUnicodeCharacter{1D5}{\ensuremath{^{u}}}
\DeclareUnicodeCharacter{02E1}{\ensuremath{^{l}}}
\DeclareUnicodeCharacter{2295}{\ensuremath{\oplus}}
\DeclareUnicodeCharacter{27FA}{\ensuremath{\Longleftrightarrow}}
\DeclareUnicodeCharacter{27F9}{\ensuremath{\Longrightarrow}}
\DeclareUnicodeCharacter{2299}{\ensuremath{\odot}}
% Additional glyphs required by v3.5 markdown sources
\DeclareUnicodeCharacter{00A7}{\S}
\DeclareUnicodeCharacter{00B7}{\ensuremath{\cdot}}
\DeclareUnicodeCharacter{00E1}{\'{a}}
\DeclareUnicodeCharacter{00E4}{\"{a}}
\DeclareUnicodeCharacter{00E8}{\`{e}}
\DeclareUnicodeCharacter{00E9}{\'{e}}
\DeclareUnicodeCharacter{00F6}{\"{o}}
\DeclareUnicodeCharacter{010D}{\v{c}}
\DeclareUnicodeCharacter{014D}{\=o}
\DeclareUnicodeCharacter{0161}{\v{s}}
\DeclareUnicodeCharacter{0302}{}
\DeclareUnicodeCharacter{0304}{}
\DeclareUnicodeCharacter{2013}{--}
\DeclareUnicodeCharacter{2016}{\ensuremath{\|}}
\DeclareUnicodeCharacter{2020}{\ensuremath{\dagger}}
\DeclareUnicodeCharacter{2102}{\ensuremath{\mathbb{C}}}
\DeclareUnicodeCharacter{210D}{\ensuremath{\mathbb{H}}}
\DeclareUnicodeCharacter{211D}{\ensuremath{\mathbb{R}}}
\DeclareUnicodeCharacter{2124}{\ensuremath{\mathbb{Z}}}
\DeclareUnicodeCharacter{2190}{\ensuremath{\leftarrow}}
\DeclareUnicodeCharacter{2227}{\ensuremath{\wedge}}
\DeclareUnicodeCharacter{2228}{\ensuremath{\vee}}
\DeclareUnicodeCharacter{2283}{\ensuremath{\supset}}
\DeclareUnicodeCharacter{22C6}{\ensuremath{\star}}
\DeclareUnicodeCharacter{22EF}{\ensuremath{\cdots}}
\DeclareUnicodeCharacter{23F3}{\texttt{[\ldots]}}
% Provide \tightlist used by pandoc-generated lists (no-op)
\providecommand{\tightlist}{\setlength{\itemsep}{0pt}\setlength{\parskip}{0pt}}
% pandoc longtable column specs: provide \real as passthrough and a
% `none` counter so that calc's \real{x} expansion path works.
\newcounter{none}
\providecommand{\real}[1]{#1}

% LISTINGS
\lstset{
    basicstyle=\small\ttfamily,
    breaklines=true, frame=single,
    keepspaces=true, showstringspaces=false,
    aboveskip=0.8em, belowskip=0.8em
}

% HEADER/FOOTER
\setlength{\headheight}{14pt}
\pagestyle{fancy}
\fancyhf{}
\fancyhead[R]{\thepage}
\renewcommand{\headrulewidth}{0.2pt}

% HYPERREF
\hypersetup{
    colorlinks=true,
    linkcolor=blue, citecolor=blue, urlcolor=blue,
    pdfauthor={Brieuc de La Fourniere}
}

% SPACING
\setstretch{1.2}
\setlength{\parskip}{0.4em}
\setlength{\parindent}{0pt}

% TITLE FORMATTING
\usepackage{titling}
\pretitle{\LARGE\bfseries}
\posttitle{\vspace{-0.4em}}
\preauthor{}
\postauthor{}
\predate{}
\postdate{}
\setlength{\droptitle}{-2.0em}

% CUSTOM COMMANDS — GIFT notation
\newcommand{\E}{\mathrm{E}}
\newcommand{\Gtwo}{\mathrm{G}_2}
\newcommand{\Kseven}{K_7}
\newcommand{\AdS}{\mathrm{AdS}}
\newcommand{\SU}{\mathrm{SU}}
\newcommand{\SO}{\mathrm{SO}}
\newcommand{\U}{\mathrm{U}}
\newcommand{\Sp}{\mathrm{Sp}}
\newcommand{\PSL}{\mathrm{PSL}}
\newcommand{\SM}{\mathrm{SM}}
\newcommand{\Tr}{\mathrm{Tr}}
\newcommand{\rank}{\mathrm{rank}}
\newcommand{\Vol}{\mathrm{Vol}}
\newcommand{\Hol}{\mathrm{Hol}}
\newcommand{\Ric}{\mathrm{Ric}}
\newcommand{\Rm}{\mathrm{Rm}}
\newcommand{\Scal}{\mathrm{Scal}}
\DeclareMathOperator{\diag}{diag}
\DeclareMathOperator{\cond}{cond}
\newcommand{\GIFT}{\textrm{GIFT}}
\newcommand{\CP}{\mathrm{CP}}
\newcommand{\CKM}{\mathrm{CKM}}
\newcommand{\PMNS}{\mathrm{PMNS}}
\newcommand{\EW}{\mathrm{EW}}
\newcommand{\Pl}{\mathrm{Pl}}
\newcommand{\DE}{\mathrm{DE}}
\newcommand{\DM}{\mathrm{DM}}
\newcommand{\KK}{\mathrm{KK}}
\newcommand{\NK}{\mathrm{NK}}
\newcommand{\TCS}{\mathrm{TCS}}
\newcommand{\PINN}{\mathrm{PINN}}
\newcommand{\btwo}{b_2}
\newcommand{\bthree}{b_3}
\newcommand{\typeI}{\textsc{I}}
\newcommand{\typeII}{\textsc{II}}
\newcommand{\typeIII}{\textsc{III}}
\newcommand{\typeIV}{\textsc{IV}}
\newcommand{\proven}{\textsc{Proven}}

% PDF bookmark safety
\pdfstringdefDisableCommands{%
  \def\textsubscript#1{#1}%
  \def\textsuperscript#1{#1}%
  \def\CP{CP}%
  \def\Gtwo{G2}%
  \def\Kseven{K7}%
  \def\GIFT{GIFT}%
  \def\E{E}%
  \def\SM{SM}%
}

% THEOREM ENVIRONMENTS
\newtheorem{theorem}{Theorem}[section]
\newtheorem{proposition}[theorem]{Proposition}
\newtheorem{lemma}[theorem]{Lemma}
\newtheorem{corollary}[theorem]{Corollary}
\theoremstyle{definition}
\newtheorem{definition}[theorem]{Definition}
\theoremstyle{remark}
\newtheorem{remark}[theorem]{Remark}
 after \documentclass

% ENCODING & FONTS
\usepackage[utf8]{inputenc}
\usepackage[T1]{fontenc}
\usepackage{lmodern}

% PAGE LAYOUT (Golden Ratio)
\usepackage[margin=1.618cm, top=2.618cm, bottom=2.618cm]{geometry}

% ESSENTIAL PACKAGES
\usepackage{float}
\usepackage{caption}
\usepackage{subcaption}
\usepackage{setspace}
\usepackage{fancyhdr}
\usepackage{xcolor}
\usepackage{hyperref}
\usepackage{csquotes}
\usepackage{amsmath}
\usepackage{amssymb}
\usepackage{amsthm}
\usepackage{mathtools}
\usepackage{booktabs}
\usepackage{longtable}
\usepackage{array}
\usepackage{multirow}
\usepackage{tikz}
\usepackage{graphicx}
\usepackage{listings}
\usepackage{enumitem}
% Unicode fallbacks (pandoc artifacts that survive markdown→LaTeX)
\DeclareUnicodeCharacter{00B0}{\ensuremath{^\circ}}
\DeclareUnicodeCharacter{221A}{\ensuremath{\sqrt{\,}}}
\DeclareUnicodeCharacter{2713}{\ensuremath{\checkmark}}
\DeclareUnicodeCharacter{2717}{\ensuremath{\times}}
\DeclareUnicodeCharacter{2070}{\ensuremath{^{0}}}
\DeclareUnicodeCharacter{00B9}{\ensuremath{^{1}}}
\DeclareUnicodeCharacter{00B2}{\ensuremath{^{2}}}
\DeclareUnicodeCharacter{00B3}{\ensuremath{^{3}}}
\DeclareUnicodeCharacter{2074}{\ensuremath{^{4}}}
\DeclareUnicodeCharacter{2075}{\ensuremath{^{5}}}
\DeclareUnicodeCharacter{2076}{\ensuremath{^{6}}}
\DeclareUnicodeCharacter{2077}{\ensuremath{^{7}}}
\DeclareUnicodeCharacter{2078}{\ensuremath{^{8}}}
\DeclareUnicodeCharacter{2079}{\ensuremath{^{9}}}
\DeclareUnicodeCharacter{2071}{\ensuremath{^{i}}}
\DeclareUnicodeCharacter{2080}{\ensuremath{_{0}}}
\DeclareUnicodeCharacter{2081}{\ensuremath{_{1}}}
\DeclareUnicodeCharacter{2082}{\ensuremath{_{2}}}
\DeclareUnicodeCharacter{2083}{\ensuremath{_{3}}}
\DeclareUnicodeCharacter{2084}{\ensuremath{_{4}}}
\DeclareUnicodeCharacter{2085}{\ensuremath{_{5}}}
\DeclareUnicodeCharacter{2086}{\ensuremath{_{6}}}
\DeclareUnicodeCharacter{2087}{\ensuremath{_{7}}}
\DeclareUnicodeCharacter{2088}{\ensuremath{_{8}}}
\DeclareUnicodeCharacter{2089}{\ensuremath{_{9}}}
\DeclareUnicodeCharacter{208A}{\ensuremath{_{+}}}
\DeclareUnicodeCharacter{208B}{\ensuremath{_{-}}}
\DeclareUnicodeCharacter{2192}{\ensuremath{\to}}
\DeclareUnicodeCharacter{2502}{|}
\DeclareUnicodeCharacter{25BC}{\ensuremath{\blacktriangledown}}
\DeclareUnicodeCharacter{2500}{\ensuremath{-}}
\DeclareUnicodeCharacter{2514}{\ensuremath{\llcorner}}
\DeclareUnicodeCharacter{251C}{\ensuremath{\vdash}}
\DeclareUnicodeCharacter{2026}{\ldots}
\DeclareUnicodeCharacter{2208}{\ensuremath{\in}}
\DeclareUnicodeCharacter{2229}{\ensuremath{\cap}}
\DeclareUnicodeCharacter{222A}{\ensuremath{\cup}}
\DeclareUnicodeCharacter{2295}{\ensuremath{\oplus}}
\DeclareUnicodeCharacter{2297}{\ensuremath{\otimes}}
\DeclareUnicodeCharacter{2245}{\ensuremath{\cong}}
\DeclareUnicodeCharacter{2248}{\ensuremath{\approx}}
\DeclareUnicodeCharacter{2260}{\ensuremath{\neq}}
\DeclareUnicodeCharacter{2261}{\ensuremath{\equiv}}
\DeclareUnicodeCharacter{2264}{\ensuremath{\leq}}
\DeclareUnicodeCharacter{2265}{\ensuremath{\geq}}
\DeclareUnicodeCharacter{00D7}{\ensuremath{\times}}
\DeclareUnicodeCharacter{00B1}{\ensuremath{\pm}}
\DeclareUnicodeCharacter{2225}{\ensuremath{\|}}
% Greek letters used in body text (pandoc passes them through unchanged)
\DeclareUnicodeCharacter{03B1}{\ensuremath{\alpha}}
\DeclareUnicodeCharacter{03B2}{\ensuremath{\beta}}
\DeclareUnicodeCharacter{03B3}{\ensuremath{\gamma}}
\DeclareUnicodeCharacter{03B4}{\ensuremath{\delta}}
\DeclareUnicodeCharacter{03B5}{\ensuremath{\varepsilon}}
\DeclareUnicodeCharacter{03B6}{\ensuremath{\zeta}}
\DeclareUnicodeCharacter{03B7}{\ensuremath{\eta}}
\DeclareUnicodeCharacter{03B8}{\ensuremath{\theta}}
\DeclareUnicodeCharacter{03B9}{\ensuremath{\iota}}
\DeclareUnicodeCharacter{03BA}{\ensuremath{\kappa}}
\DeclareUnicodeCharacter{03BB}{\ensuremath{\lambda}}
\DeclareUnicodeCharacter{03BC}{\ensuremath{\mu}}
\DeclareUnicodeCharacter{03BD}{\ensuremath{\nu}}
\DeclareUnicodeCharacter{03BE}{\ensuremath{\xi}}
\DeclareUnicodeCharacter{03C0}{\ensuremath{\pi}}
\DeclareUnicodeCharacter{03C1}{\ensuremath{\rho}}
\DeclareUnicodeCharacter{03C3}{\ensuremath{\sigma}}
\DeclareUnicodeCharacter{03C4}{\ensuremath{\tau}}
\DeclareUnicodeCharacter{03C5}{\ensuremath{\upsilon}}
\DeclareUnicodeCharacter{03C6}{\ensuremath{\varphi}}
\DeclareUnicodeCharacter{03C7}{\ensuremath{\chi}}
\DeclareUnicodeCharacter{03C8}{\ensuremath{\psi}}
\DeclareUnicodeCharacter{03C9}{\ensuremath{\omega}}
\DeclareUnicodeCharacter{0394}{\ensuremath{\Delta}}
\DeclareUnicodeCharacter{0398}{\ensuremath{\Theta}}
\DeclareUnicodeCharacter{039B}{\ensuremath{\Lambda}}
\DeclareUnicodeCharacter{03A0}{\ensuremath{\Pi}}
\DeclareUnicodeCharacter{03A3}{\ensuremath{\Sigma}}
\DeclareUnicodeCharacter{03A6}{\ensuremath{\Phi}}
\DeclareUnicodeCharacter{03A8}{\ensuremath{\Psi}}
\DeclareUnicodeCharacter{03A9}{\ensuremath{\Omega}}
\DeclareUnicodeCharacter{2200}{\ensuremath{\forall}}
\DeclareUnicodeCharacter{2203}{\ensuremath{\exists}}
\DeclareUnicodeCharacter{2207}{\ensuremath{\nabla}}
\DeclareUnicodeCharacter{2202}{\ensuremath{\partial}}
\DeclareUnicodeCharacter{2218}{\ensuremath{\circ}}
\DeclareUnicodeCharacter{2032}{\ensuremath{\prime}}
\DeclareUnicodeCharacter{2113}{\ensuremath{\ell}}
\DeclareUnicodeCharacter{210F}{\ensuremath{\hbar}}
\DeclareUnicodeCharacter{27E8}{\ensuremath{\langle}}
\DeclareUnicodeCharacter{27E9}{\ensuremath{\rangle}}
\DeclareUnicodeCharacter{2308}{\ensuremath{\lceil}}
\DeclareUnicodeCharacter{2309}{\ensuremath{\rceil}}
\DeclareUnicodeCharacter{230A}{\ensuremath{\lfloor}}
\DeclareUnicodeCharacter{230B}{\ensuremath{\rfloor}}
\DeclareUnicodeCharacter{222B}{\ensuremath{\int}}
\DeclareUnicodeCharacter{2211}{\ensuremath{\sum}}
\DeclareUnicodeCharacter{220F}{\ensuremath{\prod}}
\DeclareUnicodeCharacter{221E}{\ensuremath{\infty}}
\DeclareUnicodeCharacter{00B5}{\ensuremath{\mu}}
\DeclareUnicodeCharacter{2236}{\ensuremath{\colon}}
\DeclareUnicodeCharacter{2192}{\ensuremath{\to}}
\DeclareUnicodeCharacter{21A6}{\ensuremath{\mapsto}}
\DeclareUnicodeCharacter{21D2}{\ensuremath{\Rightarrow}}
\DeclareUnicodeCharacter{21D4}{\ensuremath{\Leftrightarrow}}
\DeclareUnicodeCharacter{2282}{\ensuremath{\subset}}
\DeclareUnicodeCharacter{2286}{\ensuremath{\subseteq}}
\DeclareUnicodeCharacter{2287}{\ensuremath{\supseteq}}
\DeclareUnicodeCharacter{222B}{\ensuremath{\int}}
\DeclareUnicodeCharacter{2205}{\ensuremath{\emptyset}}
\DeclareUnicodeCharacter{2299}{\ensuremath{\odot}}
\DeclareUnicodeCharacter{2220}{\ensuremath{\angle}}
\DeclareUnicodeCharacter{2135}{\ensuremath{\aleph}}
\DeclareUnicodeCharacter{207B}{\ensuremath{^{-}}}
\DeclareUnicodeCharacter{207A}{\ensuremath{^{+}}}
\DeclareUnicodeCharacter{220E}{\ensuremath{\blacksquare}}
\DeclareUnicodeCharacter{2061}{}
\DeclareUnicodeCharacter{2062}{}
\DeclareUnicodeCharacter{2063}{}
\DeclareUnicodeCharacter{2009}{\,}
\DeclareUnicodeCharacter{200A}{\,}
\DeclareUnicodeCharacter{200B}{}
\DeclareUnicodeCharacter{2212}{\ensuremath{-}}
\DeclareUnicodeCharacter{2194}{\ensuremath{\leftrightarrow}}
\DeclareUnicodeCharacter{2099}{\ensuremath{_{n}}}
\DeclareUnicodeCharacter{2098}{\ensuremath{_{m}}}
\DeclareUnicodeCharacter{2096}{\ensuremath{_{k}}}
\DeclareUnicodeCharacter{2095}{\ensuremath{_{h}}}
\DeclareUnicodeCharacter{2090}{\ensuremath{_{a}}}
\DeclareUnicodeCharacter{2091}{\ensuremath{_{e}}}
\DeclareUnicodeCharacter{2093}{\ensuremath{_{x}}}
\DeclareUnicodeCharacter{1D62}{\ensuremath{_{i}}}
\DeclareUnicodeCharacter{2C7C}{\ensuremath{_{j}}}
\DeclareUnicodeCharacter{2C7B}{\ensuremath{^{e}}}
\DeclareUnicodeCharacter{1D52}{\ensuremath{^{o}}}
\DeclareUnicodeCharacter{1D5}{\ensuremath{^{u}}}
\DeclareUnicodeCharacter{02E1}{\ensuremath{^{l}}}
\DeclareUnicodeCharacter{2295}{\ensuremath{\oplus}}
\DeclareUnicodeCharacter{27FA}{\ensuremath{\Longleftrightarrow}}
\DeclareUnicodeCharacter{27F9}{\ensuremath{\Longrightarrow}}
\DeclareUnicodeCharacter{2299}{\ensuremath{\odot}}
% Additional glyphs required by v3.5 markdown sources
\DeclareUnicodeCharacter{00A7}{\S}
\DeclareUnicodeCharacter{00B7}{\ensuremath{\cdot}}
\DeclareUnicodeCharacter{00E1}{\'{a}}
\DeclareUnicodeCharacter{00E4}{\"{a}}
\DeclareUnicodeCharacter{00E8}{\`{e}}
\DeclareUnicodeCharacter{00E9}{\'{e}}
\DeclareUnicodeCharacter{00F6}{\"{o}}
\DeclareUnicodeCharacter{010D}{\v{c}}
\DeclareUnicodeCharacter{014D}{\=o}
\DeclareUnicodeCharacter{0161}{\v{s}}
\DeclareUnicodeCharacter{0302}{}
\DeclareUnicodeCharacter{0304}{}
\DeclareUnicodeCharacter{2013}{--}
\DeclareUnicodeCharacter{2016}{\ensuremath{\|}}
\DeclareUnicodeCharacter{2020}{\ensuremath{\dagger}}
\DeclareUnicodeCharacter{2102}{\ensuremath{\mathbb{C}}}
\DeclareUnicodeCharacter{210D}{\ensuremath{\mathbb{H}}}
\DeclareUnicodeCharacter{211D}{\ensuremath{\mathbb{R}}}
\DeclareUnicodeCharacter{2124}{\ensuremath{\mathbb{Z}}}
\DeclareUnicodeCharacter{2190}{\ensuremath{\leftarrow}}
\DeclareUnicodeCharacter{2227}{\ensuremath{\wedge}}
\DeclareUnicodeCharacter{2228}{\ensuremath{\vee}}
\DeclareUnicodeCharacter{2283}{\ensuremath{\supset}}
\DeclareUnicodeCharacter{22C6}{\ensuremath{\star}}
\DeclareUnicodeCharacter{22EF}{\ensuremath{\cdots}}
\DeclareUnicodeCharacter{23F3}{\texttt{[\ldots]}}
% Provide \tightlist used by pandoc-generated lists (no-op)
\providecommand{\tightlist}{\setlength{\itemsep}{0pt}\setlength{\parskip}{0pt}}
% pandoc longtable column specs: provide \real as passthrough and a
% `none` counter so that calc's \real{x} expansion path works.
\newcounter{none}
\providecommand{\real}[1]{#1}

% LISTINGS
\lstset{
    basicstyle=\small\ttfamily,
    breaklines=true, frame=single,
    keepspaces=true, showstringspaces=false,
    aboveskip=0.8em, belowskip=0.8em
}

% HEADER/FOOTER
\setlength{\headheight}{14pt}
\pagestyle{fancy}
\fancyhf{}
\fancyhead[R]{\thepage}
\renewcommand{\headrulewidth}{0.2pt}

% HYPERREF
\hypersetup{
    colorlinks=true,
    linkcolor=blue, citecolor=blue, urlcolor=blue,
    pdfauthor={Brieuc de La Fourniere}
}

% SPACING
\setstretch{1.2}
\setlength{\parskip}{0.4em}
\setlength{\parindent}{0pt}

% TITLE FORMATTING
\usepackage{titling}
\pretitle{\LARGE\bfseries}
\posttitle{\vspace{-0.4em}}
\preauthor{}
\postauthor{}
\predate{}
\postdate{}
\setlength{\droptitle}{-2.0em}

% CUSTOM COMMANDS — GIFT notation
\newcommand{\E}{\mathrm{E}}
\newcommand{\Gtwo}{\mathrm{G}_2}
\newcommand{\Kseven}{K_7}
\newcommand{\AdS}{\mathrm{AdS}}
\newcommand{\SU}{\mathrm{SU}}
\newcommand{\SO}{\mathrm{SO}}
\newcommand{\U}{\mathrm{U}}
\newcommand{\Sp}{\mathrm{Sp}}
\newcommand{\PSL}{\mathrm{PSL}}
\newcommand{\SM}{\mathrm{SM}}
\newcommand{\Tr}{\mathrm{Tr}}
\newcommand{\rank}{\mathrm{rank}}
\newcommand{\Vol}{\mathrm{Vol}}
\newcommand{\Hol}{\mathrm{Hol}}
\newcommand{\Ric}{\mathrm{Ric}}
\newcommand{\Rm}{\mathrm{Rm}}
\newcommand{\Scal}{\mathrm{Scal}}
\DeclareMathOperator{\diag}{diag}
\DeclareMathOperator{\cond}{cond}
\newcommand{\GIFT}{\textrm{GIFT}}
\newcommand{\CP}{\mathrm{CP}}
\newcommand{\CKM}{\mathrm{CKM}}
\newcommand{\PMNS}{\mathrm{PMNS}}
\newcommand{\EW}{\mathrm{EW}}
\newcommand{\Pl}{\mathrm{Pl}}
\newcommand{\DE}{\mathrm{DE}}
\newcommand{\DM}{\mathrm{DM}}
\newcommand{\KK}{\mathrm{KK}}
\newcommand{\NK}{\mathrm{NK}}
\newcommand{\TCS}{\mathrm{TCS}}
\newcommand{\PINN}{\mathrm{PINN}}
\newcommand{\btwo}{b_2}
\newcommand{\bthree}{b_3}
\newcommand{\typeI}{\textsc{I}}
\newcommand{\typeII}{\textsc{II}}
\newcommand{\typeIII}{\textsc{III}}
\newcommand{\typeIV}{\textsc{IV}}
\newcommand{\proven}{\textsc{Proven}}

% PDF bookmark safety
\pdfstringdefDisableCommands{%
  \def\textsubscript#1{#1}%
  \def\textsuperscript#1{#1}%
  \def\CP{CP}%
  \def\Gtwo{G2}%
  \def\Kseven{K7}%
  \def\GIFT{GIFT}%
  \def\E{E}%
  \def\SM{SM}%
}

% THEOREM ENVIRONMENTS
\newtheorem{theorem}{Theorem}[section]
\newtheorem{proposition}[theorem]{Proposition}
\newtheorem{lemma}[theorem]{Lemma}
\newtheorem{corollary}[theorem]{Corollary}
\theoremstyle{definition}
\newtheorem{definition}[theorem]{Definition}
\theoremstyle{remark}
\newtheorem{remark}[theorem]{Remark}
 after \documentclass

% ENCODING & FONTS
\usepackage[utf8]{inputenc}
\usepackage[T1]{fontenc}
\usepackage{lmodern}

% PAGE LAYOUT (Golden Ratio)
\usepackage[margin=1.618cm, top=2.618cm, bottom=2.618cm]{geometry}

% ESSENTIAL PACKAGES
\usepackage{float}
\usepackage{caption}
\usepackage{subcaption}
\usepackage{setspace}
\usepackage{fancyhdr}
\usepackage{xcolor}
\usepackage{hyperref}
\usepackage{csquotes}
\usepackage{amsmath}
\usepackage{amssymb}
\usepackage{amsthm}
\usepackage{mathtools}
\usepackage{booktabs}
\usepackage{longtable}
\usepackage{array}
\usepackage{multirow}
\usepackage{tikz}
\usepackage{graphicx}
\usepackage{listings}
\usepackage{enumitem}
% Unicode fallbacks (pandoc artifacts that survive markdown→LaTeX)
\DeclareUnicodeCharacter{00B0}{\ensuremath{^\circ}}
\DeclareUnicodeCharacter{221A}{\ensuremath{\sqrt{\,}}}
\DeclareUnicodeCharacter{2713}{\ensuremath{\checkmark}}
\DeclareUnicodeCharacter{2717}{\ensuremath{\times}}
\DeclareUnicodeCharacter{2070}{\ensuremath{^{0}}}
\DeclareUnicodeCharacter{00B9}{\ensuremath{^{1}}}
\DeclareUnicodeCharacter{00B2}{\ensuremath{^{2}}}
\DeclareUnicodeCharacter{00B3}{\ensuremath{^{3}}}
\DeclareUnicodeCharacter{2074}{\ensuremath{^{4}}}
\DeclareUnicodeCharacter{2075}{\ensuremath{^{5}}}
\DeclareUnicodeCharacter{2076}{\ensuremath{^{6}}}
\DeclareUnicodeCharacter{2077}{\ensuremath{^{7}}}
\DeclareUnicodeCharacter{2078}{\ensuremath{^{8}}}
\DeclareUnicodeCharacter{2079}{\ensuremath{^{9}}}
\DeclareUnicodeCharacter{2071}{\ensuremath{^{i}}}
\DeclareUnicodeCharacter{2080}{\ensuremath{_{0}}}
\DeclareUnicodeCharacter{2081}{\ensuremath{_{1}}}
\DeclareUnicodeCharacter{2082}{\ensuremath{_{2}}}
\DeclareUnicodeCharacter{2083}{\ensuremath{_{3}}}
\DeclareUnicodeCharacter{2084}{\ensuremath{_{4}}}
\DeclareUnicodeCharacter{2085}{\ensuremath{_{5}}}
\DeclareUnicodeCharacter{2086}{\ensuremath{_{6}}}
\DeclareUnicodeCharacter{2087}{\ensuremath{_{7}}}
\DeclareUnicodeCharacter{2088}{\ensuremath{_{8}}}
\DeclareUnicodeCharacter{2089}{\ensuremath{_{9}}}
\DeclareUnicodeCharacter{208A}{\ensuremath{_{+}}}
\DeclareUnicodeCharacter{208B}{\ensuremath{_{-}}}
\DeclareUnicodeCharacter{2192}{\ensuremath{\to}}
\DeclareUnicodeCharacter{2502}{|}
\DeclareUnicodeCharacter{25BC}{\ensuremath{\blacktriangledown}}
\DeclareUnicodeCharacter{2500}{\ensuremath{-}}
\DeclareUnicodeCharacter{2514}{\ensuremath{\llcorner}}
\DeclareUnicodeCharacter{251C}{\ensuremath{\vdash}}
\DeclareUnicodeCharacter{2026}{\ldots}
\DeclareUnicodeCharacter{2208}{\ensuremath{\in}}
\DeclareUnicodeCharacter{2229}{\ensuremath{\cap}}
\DeclareUnicodeCharacter{222A}{\ensuremath{\cup}}
\DeclareUnicodeCharacter{2295}{\ensuremath{\oplus}}
\DeclareUnicodeCharacter{2297}{\ensuremath{\otimes}}
\DeclareUnicodeCharacter{2245}{\ensuremath{\cong}}
\DeclareUnicodeCharacter{2248}{\ensuremath{\approx}}
\DeclareUnicodeCharacter{2260}{\ensuremath{\neq}}
\DeclareUnicodeCharacter{2261}{\ensuremath{\equiv}}
\DeclareUnicodeCharacter{2264}{\ensuremath{\leq}}
\DeclareUnicodeCharacter{2265}{\ensuremath{\geq}}
\DeclareUnicodeCharacter{00D7}{\ensuremath{\times}}
\DeclareUnicodeCharacter{00B1}{\ensuremath{\pm}}
\DeclareUnicodeCharacter{2225}{\ensuremath{\|}}
% Greek letters used in body text (pandoc passes them through unchanged)
\DeclareUnicodeCharacter{03B1}{\ensuremath{\alpha}}
\DeclareUnicodeCharacter{03B2}{\ensuremath{\beta}}
\DeclareUnicodeCharacter{03B3}{\ensuremath{\gamma}}
\DeclareUnicodeCharacter{03B4}{\ensuremath{\delta}}
\DeclareUnicodeCharacter{03B5}{\ensuremath{\varepsilon}}
\DeclareUnicodeCharacter{03B6}{\ensuremath{\zeta}}
\DeclareUnicodeCharacter{03B7}{\ensuremath{\eta}}
\DeclareUnicodeCharacter{03B8}{\ensuremath{\theta}}
\DeclareUnicodeCharacter{03B9}{\ensuremath{\iota}}
\DeclareUnicodeCharacter{03BA}{\ensuremath{\kappa}}
\DeclareUnicodeCharacter{03BB}{\ensuremath{\lambda}}
\DeclareUnicodeCharacter{03BC}{\ensuremath{\mu}}
\DeclareUnicodeCharacter{03BD}{\ensuremath{\nu}}
\DeclareUnicodeCharacter{03BE}{\ensuremath{\xi}}
\DeclareUnicodeCharacter{03C0}{\ensuremath{\pi}}
\DeclareUnicodeCharacter{03C1}{\ensuremath{\rho}}
\DeclareUnicodeCharacter{03C3}{\ensuremath{\sigma}}
\DeclareUnicodeCharacter{03C4}{\ensuremath{\tau}}
\DeclareUnicodeCharacter{03C5}{\ensuremath{\upsilon}}
\DeclareUnicodeCharacter{03C6}{\ensuremath{\varphi}}
\DeclareUnicodeCharacter{03C7}{\ensuremath{\chi}}
\DeclareUnicodeCharacter{03C8}{\ensuremath{\psi}}
\DeclareUnicodeCharacter{03C9}{\ensuremath{\omega}}
\DeclareUnicodeCharacter{0394}{\ensuremath{\Delta}}
\DeclareUnicodeCharacter{0398}{\ensuremath{\Theta}}
\DeclareUnicodeCharacter{039B}{\ensuremath{\Lambda}}
\DeclareUnicodeCharacter{03A0}{\ensuremath{\Pi}}
\DeclareUnicodeCharacter{03A3}{\ensuremath{\Sigma}}
\DeclareUnicodeCharacter{03A6}{\ensuremath{\Phi}}
\DeclareUnicodeCharacter{03A8}{\ensuremath{\Psi}}
\DeclareUnicodeCharacter{03A9}{\ensuremath{\Omega}}
\DeclareUnicodeCharacter{2200}{\ensuremath{\forall}}
\DeclareUnicodeCharacter{2203}{\ensuremath{\exists}}
\DeclareUnicodeCharacter{2207}{\ensuremath{\nabla}}
\DeclareUnicodeCharacter{2202}{\ensuremath{\partial}}
\DeclareUnicodeCharacter{2218}{\ensuremath{\circ}}
\DeclareUnicodeCharacter{2032}{\ensuremath{\prime}}
\DeclareUnicodeCharacter{2113}{\ensuremath{\ell}}
\DeclareUnicodeCharacter{210F}{\ensuremath{\hbar}}
\DeclareUnicodeCharacter{27E8}{\ensuremath{\langle}}
\DeclareUnicodeCharacter{27E9}{\ensuremath{\rangle}}
\DeclareUnicodeCharacter{2308}{\ensuremath{\lceil}}
\DeclareUnicodeCharacter{2309}{\ensuremath{\rceil}}
\DeclareUnicodeCharacter{230A}{\ensuremath{\lfloor}}
\DeclareUnicodeCharacter{230B}{\ensuremath{\rfloor}}
\DeclareUnicodeCharacter{222B}{\ensuremath{\int}}
\DeclareUnicodeCharacter{2211}{\ensuremath{\sum}}
\DeclareUnicodeCharacter{220F}{\ensuremath{\prod}}
\DeclareUnicodeCharacter{221E}{\ensuremath{\infty}}
\DeclareUnicodeCharacter{00B5}{\ensuremath{\mu}}
\DeclareUnicodeCharacter{2236}{\ensuremath{\colon}}
\DeclareUnicodeCharacter{2192}{\ensuremath{\to}}
\DeclareUnicodeCharacter{21A6}{\ensuremath{\mapsto}}
\DeclareUnicodeCharacter{21D2}{\ensuremath{\Rightarrow}}
\DeclareUnicodeCharacter{21D4}{\ensuremath{\Leftrightarrow}}
\DeclareUnicodeCharacter{2282}{\ensuremath{\subset}}
\DeclareUnicodeCharacter{2286}{\ensuremath{\subseteq}}
\DeclareUnicodeCharacter{2287}{\ensuremath{\supseteq}}
\DeclareUnicodeCharacter{222B}{\ensuremath{\int}}
\DeclareUnicodeCharacter{2205}{\ensuremath{\emptyset}}
\DeclareUnicodeCharacter{2299}{\ensuremath{\odot}}
\DeclareUnicodeCharacter{2220}{\ensuremath{\angle}}
\DeclareUnicodeCharacter{2135}{\ensuremath{\aleph}}
\DeclareUnicodeCharacter{207B}{\ensuremath{^{-}}}
\DeclareUnicodeCharacter{207A}{\ensuremath{^{+}}}
\DeclareUnicodeCharacter{220E}{\ensuremath{\blacksquare}}
\DeclareUnicodeCharacter{2061}{}
\DeclareUnicodeCharacter{2062}{}
\DeclareUnicodeCharacter{2063}{}
\DeclareUnicodeCharacter{2009}{\,}
\DeclareUnicodeCharacter{200A}{\,}
\DeclareUnicodeCharacter{200B}{}
\DeclareUnicodeCharacter{2212}{\ensuremath{-}}
\DeclareUnicodeCharacter{2194}{\ensuremath{\leftrightarrow}}
\DeclareUnicodeCharacter{2099}{\ensuremath{_{n}}}
\DeclareUnicodeCharacter{2098}{\ensuremath{_{m}}}
\DeclareUnicodeCharacter{2096}{\ensuremath{_{k}}}
\DeclareUnicodeCharacter{2095}{\ensuremath{_{h}}}
\DeclareUnicodeCharacter{2090}{\ensuremath{_{a}}}
\DeclareUnicodeCharacter{2091}{\ensuremath{_{e}}}
\DeclareUnicodeCharacter{2093}{\ensuremath{_{x}}}
\DeclareUnicodeCharacter{1D62}{\ensuremath{_{i}}}
\DeclareUnicodeCharacter{2C7C}{\ensuremath{_{j}}}
\DeclareUnicodeCharacter{2C7B}{\ensuremath{^{e}}}
\DeclareUnicodeCharacter{1D52}{\ensuremath{^{o}}}
\DeclareUnicodeCharacter{1D5}{\ensuremath{^{u}}}
\DeclareUnicodeCharacter{02E1}{\ensuremath{^{l}}}
\DeclareUnicodeCharacter{2295}{\ensuremath{\oplus}}
\DeclareUnicodeCharacter{27FA}{\ensuremath{\Longleftrightarrow}}
\DeclareUnicodeCharacter{27F9}{\ensuremath{\Longrightarrow}}
\DeclareUnicodeCharacter{2299}{\ensuremath{\odot}}
% Additional glyphs required by v3.5 markdown sources
\DeclareUnicodeCharacter{00A7}{\S}
\DeclareUnicodeCharacter{00B7}{\ensuremath{\cdot}}
\DeclareUnicodeCharacter{00E1}{\'{a}}
\DeclareUnicodeCharacter{00E4}{\"{a}}
\DeclareUnicodeCharacter{00E8}{\`{e}}
\DeclareUnicodeCharacter{00E9}{\'{e}}
\DeclareUnicodeCharacter{00F6}{\"{o}}
\DeclareUnicodeCharacter{010D}{\v{c}}
\DeclareUnicodeCharacter{014D}{\=o}
\DeclareUnicodeCharacter{0161}{\v{s}}
\DeclareUnicodeCharacter{0302}{}
\DeclareUnicodeCharacter{0304}{}
\DeclareUnicodeCharacter{2013}{--}
\DeclareUnicodeCharacter{2016}{\ensuremath{\|}}
\DeclareUnicodeCharacter{2020}{\ensuremath{\dagger}}
\DeclareUnicodeCharacter{2102}{\ensuremath{\mathbb{C}}}
\DeclareUnicodeCharacter{210D}{\ensuremath{\mathbb{H}}}
\DeclareUnicodeCharacter{211D}{\ensuremath{\mathbb{R}}}
\DeclareUnicodeCharacter{2124}{\ensuremath{\mathbb{Z}}}
\DeclareUnicodeCharacter{2190}{\ensuremath{\leftarrow}}
\DeclareUnicodeCharacter{2227}{\ensuremath{\wedge}}
\DeclareUnicodeCharacter{2228}{\ensuremath{\vee}}
\DeclareUnicodeCharacter{2283}{\ensuremath{\supset}}
\DeclareUnicodeCharacter{22C6}{\ensuremath{\star}}
\DeclareUnicodeCharacter{22EF}{\ensuremath{\cdots}}
\DeclareUnicodeCharacter{23F3}{\texttt{[\ldots]}}
% Provide \tightlist used by pandoc-generated lists (no-op)
\providecommand{\tightlist}{\setlength{\itemsep}{0pt}\setlength{\parskip}{0pt}}
% pandoc longtable column specs: provide \real as passthrough and a
% `none` counter so that calc's \real{x} expansion path works.
\newcounter{none}
\providecommand{\real}[1]{#1}

% LISTINGS
\lstset{
    basicstyle=\small\ttfamily,
    breaklines=true, frame=single,
    keepspaces=true, showstringspaces=false,
    aboveskip=0.8em, belowskip=0.8em
}

% HEADER/FOOTER
\setlength{\headheight}{14pt}
\pagestyle{fancy}
\fancyhf{}
\fancyhead[R]{\thepage}
\renewcommand{\headrulewidth}{0.2pt}

% HYPERREF
\hypersetup{
    colorlinks=true,
    linkcolor=blue, citecolor=blue, urlcolor=blue,
    pdfauthor={Brieuc de La Fourniere}
}

% SPACING
\setstretch{1.2}
\setlength{\parskip}{0.4em}
\setlength{\parindent}{0pt}

% TITLE FORMATTING
\usepackage{titling}
\pretitle{\LARGE\bfseries}
\posttitle{\vspace{-0.4em}}
\preauthor{}
\postauthor{}
\predate{}
\postdate{}
\setlength{\droptitle}{-2.0em}

% CUSTOM COMMANDS — GIFT notation
\newcommand{\E}{\mathrm{E}}
\newcommand{\Gtwo}{\mathrm{G}_2}
\newcommand{\Kseven}{K_7}
\newcommand{\AdS}{\mathrm{AdS}}
\newcommand{\SU}{\mathrm{SU}}
\newcommand{\SO}{\mathrm{SO}}
\newcommand{\U}{\mathrm{U}}
\newcommand{\Sp}{\mathrm{Sp}}
\newcommand{\PSL}{\mathrm{PSL}}
\newcommand{\SM}{\mathrm{SM}}
\newcommand{\Tr}{\mathrm{Tr}}
\newcommand{\rank}{\mathrm{rank}}
\newcommand{\Vol}{\mathrm{Vol}}
\newcommand{\Hol}{\mathrm{Hol}}
\newcommand{\Ric}{\mathrm{Ric}}
\newcommand{\Rm}{\mathrm{Rm}}
\newcommand{\Scal}{\mathrm{Scal}}
\DeclareMathOperator{\diag}{diag}
\DeclareMathOperator{\cond}{cond}
\newcommand{\GIFT}{\textrm{GIFT}}
\newcommand{\CP}{\mathrm{CP}}
\newcommand{\CKM}{\mathrm{CKM}}
\newcommand{\PMNS}{\mathrm{PMNS}}
\newcommand{\EW}{\mathrm{EW}}
\newcommand{\Pl}{\mathrm{Pl}}
\newcommand{\DE}{\mathrm{DE}}
\newcommand{\DM}{\mathrm{DM}}
\newcommand{\KK}{\mathrm{KK}}
\newcommand{\NK}{\mathrm{NK}}
\newcommand{\TCS}{\mathrm{TCS}}
\newcommand{\PINN}{\mathrm{PINN}}
\newcommand{\btwo}{b_2}
\newcommand{\bthree}{b_3}
\newcommand{\typeI}{\textsc{I}}
\newcommand{\typeII}{\textsc{II}}
\newcommand{\typeIII}{\textsc{III}}
\newcommand{\typeIV}{\textsc{IV}}
\newcommand{\proven}{\textsc{Proven}}

% PDF bookmark safety
\pdfstringdefDisableCommands{%
  \def\textsubscript#1{#1}%
  \def\textsuperscript#1{#1}%
  \def\CP{CP}%
  \def\Gtwo{G2}%
  \def\Kseven{K7}%
  \def\GIFT{GIFT}%
  \def\E{E}%
  \def\SM{SM}%
}

% THEOREM ENVIRONMENTS
\newtheorem{theorem}{Theorem}[section]
\newtheorem{proposition}[theorem]{Proposition}
\newtheorem{lemma}[theorem]{Lemma}
\newtheorem{corollary}[theorem]{Corollary}
\theoremstyle{definition}
\newtheorem{definition}[theorem]{Definition}
\theoremstyle{remark}
\newtheorem{remark}[theorem]{Remark}

% Glyphs new in 3.6, kept here rather than in the preamble shared with 3.4 and 3.5.
\DeclareUnicodeCharacter{02E2}{\ensuremath{^{\mathrm{S}}}}
\DeclareUnicodeCharacter{1D37}{\ensuremath{^{\mathrm{K}}}}
\DeclareUnicodeCharacter{1D39}{\ensuremath{^{\mathrm{M}}}}
\DeclareUnicodeCharacter{1D3A}{\ensuremath{^{\mathrm{N}}}}
\DeclareUnicodeCharacter{1D3E}{\ensuremath{^{\mathrm{P}}}}
\DeclareUnicodeCharacter{1D9C}{\ensuremath{^{\mathrm{C}}}}
\DeclareUnicodeCharacter{211A}{\ensuremath{\mathbb{Q}}}
\DeclareUnicodeCharacter{226B}{\ensuremath{\gg}}
\DeclareUnicodeCharacter{22A5}{\ensuremath{\perp}}
% pandoc >= 3.2.1 emits \pandocbounded around images; provide a fallback.
% Also clamp all images to \linewidth (natural-size figures overflowed the
% right margin -- the "truncated figures" flagged by the round-2 reviews).
% Kept here rather than in gift_preamble.tex, shared with the 3.4 build.
\providecommand{\pandocbounded}[1]{#1}
\makeatletter
\def\maxwidth{\ifdim\Gin@nat@width>\linewidth\linewidth\else\Gin@nat@width\fi}
\makeatother
\setkeys{Gin}{width=\maxwidth,keepaspectratio}
% md headings carry their own literal numbering ("3.2 What ..."); suppress
% the LaTeX section counters to avoid "0.2 1. Introduction" double numbering.
\setcounter{secnumdepth}{0}
% Loosen justification slightly: long technical identifiers (made breakable
% at build time via \allowbreak) still need extra stretch to avoid overfulls.
\emergencystretch=3em
\fancyhead[L]{The $K_7$ Dictionary v3.6: Supplement S2}
\hypersetup{pdftitle={The K7 Dictionary v3.6 Supplement S2: The 33 Relations}}

\graphicspath{{../../../figures/}}

\title{%
\LARGE\textbf{Supplement S2: The 33 Relations}\\[0.3em]
\large Arithmetic, Dictionary Choices and Status%
}
\author{Brieuc de La Fourni\`ere \\ \textit{Independent researcher, Beaune, France}}
\date{November 2026 \quad\textbar\quad v3.6}

\begin{document}

\maketitle
\noindent\rule{\textwidth}{0.2pt}

\thispagestyle{fancy}

\section{S2.0 What this supplement establishes}\label{s2.0-what-this-supplement-establishes}

Earlier versions of this supplement were titled ``Complete Derivations'' and presented each relation as a theorem with a proof. Version 3.6 withdraws that framing. No geometric map is known that sends the invariants of (K₇, E₈×E₈) to the Standard-Model parameters, and no rule assigns a formula to an observable (main §3.2). What can be established for each relation is narrower, and this supplement establishes it:

\begin{enumerate}
\def\labelenumi{\arabic{enumi}.}
\tightlist
\item
  \textbf{The arithmetic.} Given the declared constants, the closed form evaluates to the stated number. The arithmetic below is exact. Each relation has a Lean statement tying its closed form to its value (\texttt{relations\_\allowbreak{}v36\_\allowbreak{}certified}, Arithmon/K7-Lean); the form of each statement (exact rational equality, real closed form, bounds, integer) is given in S3.7.
\item
  \textbf{The dictionary choices.} Which constants the formula uses, which operations beyond integer ratios, and which entries of the residue register it depends on (main §3.3: R1--R12).
\item
  \textbf{The status.} The status of each relation against the most recent data (main §4.1, generated from the dataset) is shown in the last column of each table, with its pull. The column is generated from main §4.1 (\texttt{k7\_\allowbreak{}v3\_\allowbreak{}6\_\allowbreak{}s2\_\allowbreak{}status.\allowbreak{}py}) and is never edited here, so that the status keeps one source; the status vocabulary is defined in main §4.1.
\end{enumerate}

What a relation's arithmetic does not establish: that the formula is the right one for the observable, that the structure requires it, or that the agreement with measurement is evidence for the framework (main §7). Section S2.7 records the alternative formulas found in the corpus for the same observables; their number is a measure of the vocabulary's freedom, not of robustness.

\begin{center}\rule{\linewidth}{0.5pt}\end{center}

\section{S2.1 Notation}\label{s2.1-notation}

The twenty structural constants and their status (Lie datum, topological, input, combination) are listed in main §3.1. The relations also use the following, each a combination chosen by the framework:

\begin{longtable}[]{@{}
  >{\raggedright\arraybackslash}p{(\linewidth - 6\tabcolsep) * \real{0.2222}}
  >{\raggedright\arraybackslash}p{(\linewidth - 6\tabcolsep) * \real{0.1944}}
  >{\raggedright\arraybackslash}p{(\linewidth - 6\tabcolsep) * \real{0.3056}}
  >{\raggedright\arraybackslash}p{(\linewidth - 6\tabcolsep) * \real{0.2778}}@{}}
\toprule\noalign{}
\begin{minipage}[b]{\linewidth}\raggedright
Symbol
\end{minipage} & \begin{minipage}[b]{\linewidth}\raggedright
Value
\end{minipage} & \begin{minipage}[b]{\linewidth}\raggedright
Definition
\end{minipage} & \begin{minipage}[b]{\linewidth}\raggedright
Register
\end{minipage} \\
\midrule\noalign{}
\endhead
\bottomrule\noalign{}
\endlastfoot
det(g) & 65/32 & p₂ + 1/(b₂ + dim G₂ − N\textsubscript{gen}) = det\textsubscript{num}/det\textsubscript{den} (main App. B) & \textbf{imposed} normalization (R2); a second expression, Weyl · α\textsubscript{sum}/2\textsuperscript{Weyl}, agrees with it on the line b₂ = 21 and differs elsewhere \\
κ\textsubscript{T}⁻¹ & 61 & b₃ − dim(G₂) − p₂ & integer combination without a physical name (R3); a second expression, b₃ − b₂ + Weyl, agrees with it on the line b₂ = 21 and differs elsewhere \\
φ & (1 + √5)/2 & through the McKay correspondence E₈ ↔ 2I & additional step (R9) \\
1 & 1 & written b₀ in earlier versions & used freely as an additive constant \\
2b₂ & 42 & also p₂ × N\textsubscript{gen} × dim(K₇) & combination \\
τ & 3472/891 & dim(E₈×E₈) · b₂/(dim J₃(\ensuremath{\mathbb{O}}) · H*) & used only in Type II entries (S3) \\
\end{longtable}

\textbf{Weyl = 5.} Three expressions give 5: (dim G₂ + 1)/N\textsubscript{gen}, b₂/N\textsubscript{gen} − p₂, dim G₂ − rank E₈ − 1. Earlier versions called this a ``triple identity''. It is an arithmetic coincidence among the declared integers; it fixes no further quantity.

\textbf{N\textsubscript{gen} = 3.} The integer identity (rank E₈ + N\textsubscript{gen}) · b₂ = N\textsubscript{gen} · b₃, i.e.~11 · 21 = 3 · 77, is satisfied by N\textsubscript{gen} = 3. It is a consistency check on (21, 77), not the output of an index theorem: the ``Atiyah--Singer'' expression of earlier versions, (77 − 8/3 · 21)/7 = 3, is the same arithmetic written as an index without an operator whose index it is (main §2.4).

\begin{center}\rule{\linewidth}{0.5pt}\end{center}

\section{S2.2 Gauge sector and Higgs}\label{s2.2-gauge-sector-and-higgs}

{\footnotesize\setlength{\tabcolsep}{3pt}
\begin{longtable}[]{@{}
  >{\raggedright\arraybackslash}p{(\linewidth - 8\tabcolsep) * \real{0.0870}}
  >{\raggedright\arraybackslash}p{(\linewidth - 8\tabcolsep) * \real{0.2174}}
  >{\raggedright\arraybackslash}p{(\linewidth - 8\tabcolsep) * \real{0.2391}}
  >{\raggedright\arraybackslash}p{(\linewidth - 8\tabcolsep) * \real{0.3913}}
  >{\raggedright\arraybackslash}p{(\linewidth - 8\tabcolsep) * \real{0.0652}}@{}}
\toprule\noalign{}
\begin{minipage}[b]{\linewidth}\raggedright
ID
\end{minipage} & \begin{minipage}[b]{\linewidth}\raggedright
Relation
\end{minipage} & \begin{minipage}[b]{\linewidth}\raggedright
Arithmetic
\end{minipage} & \begin{minipage}[b]{\linewidth}\raggedright
Choices declared
\end{minipage} & \begin{minipage}[b]{\linewidth}\raggedright
Status (main §4.1)
\end{minipage} \\
\midrule\noalign{}
\endhead
\bottomrule\noalign{}
\endlastfoot
A01 & sin²θ\textsubscript{W} (M\textsubscript{Z}, MS-bar) & b₂/(b₃ + dim G₂) = 21/91 = \textbf{3/13} = 0.230769 & scale and scheme (R7) & integer; \textbf{proximity}; pull −7.5σ \\
A02 & α\textsubscript{s}(M\textsubscript{Z}) & √2/(dim G₂ − p₂) = \textbf{√2/12} = 0.117851 & operation √; scale (R7) & algebraic; pull −0.2σ \\
A03 & λ\textsubscript{H} & √(dim G₂ + N\textsubscript{gen})/det\textsubscript{den} = \textbf{√17/32} = 0.128847 & operation √; det(g) (R2); compared with the tree-level m\textsubscript{H}²/(2v²) & algebraic; det(g); pull −1.3σ \\
A04 & α⁻¹(0) & (dim E₈ + rank E₈)/2 + H*/D\textsubscript{bulk} + det(g) · κ\textsubscript{T} = 128 + 9 + 65/1952 = \textbf{267489/1952} = 137.03330 & det(g) (R2); κ\textsubscript{T} (R3); scale q² = 0 (R7) & det(g), κ\textsubscript{T}; \textbf{proximity}; pull ≫ 10σ \\
\end{longtable}}

\textbf{On A04.} The three terms were written as ``algebraic'', ``bulk impedance'' and ``torsional correction'' in earlier versions. None of the three names corresponds to a computation; the correction term is the product of the imposed normalization and the unnamed combination. The relation lies 19.7 ppm below α⁻¹(0), whose uncertainty is 1.5 × 10⁻¹⁰ relative: a proximity (main §4.2).

\textbf{On A03.} A second formula for λ\textsubscript{H} is recorded in the corpus (S2.7).

\begin{center}\rule{\linewidth}{0.5pt}\end{center}

\section{S2.3 Leptons (overdetermined: R8)}\label{s2.3-leptons-overdetermined-r8}

{\footnotesize\setlength{\tabcolsep}{3pt}
\begin{longtable}[]{@{}
  >{\raggedright\arraybackslash}p{(\linewidth - 8\tabcolsep) * \real{0.0870}}
  >{\raggedright\arraybackslash}p{(\linewidth - 8\tabcolsep) * \real{0.2174}}
  >{\raggedright\arraybackslash}p{(\linewidth - 8\tabcolsep) * \real{0.2391}}
  >{\raggedright\arraybackslash}p{(\linewidth - 8\tabcolsep) * \real{0.3913}}
  >{\raggedright\arraybackslash}p{(\linewidth - 8\tabcolsep) * \real{0.0652}}@{}}
\toprule\noalign{}
\begin{minipage}[b]{\linewidth}\raggedright
ID
\end{minipage} & \begin{minipage}[b]{\linewidth}\raggedright
Relation
\end{minipage} & \begin{minipage}[b]{\linewidth}\raggedright
Arithmetic
\end{minipage} & \begin{minipage}[b]{\linewidth}\raggedright
Choices declared
\end{minipage} & \begin{minipage}[b]{\linewidth}\raggedright
Status (main §4.1)
\end{minipage} \\
\midrule\noalign{}
\endhead
\bottomrule\noalign{}
\endlastfoot
A05 & Q\textsubscript{Koide} & dim G₂/b₂ = 14/21 = \textbf{2/3} & R8 & integer; approximate (R8); pull +0.4σ \\
A06 & m\textsubscript{τ}/m\textsubscript{e} & dim K₇ + 10 · dim E₈ + 10 · H* = 7 + 2480 + 990 = \textbf{3477} & R8; the coefficient 10 is not among the twenty constants & integer; approximate (R8); \textbf{tension}; pull −2.1σ \\
A07 & m\textsubscript{μ}/m\textsubscript{e} & dim(J₃(\ensuremath{\mathbb{O}}))\textsuperscript{φ} = \textbf{27\textsuperscript{φ}} = 207.0119 & R8; φ (R9) & φ; approximate (R8); \textbf{proximity}; pull ≫ 10σ \\
A22 & m\textsubscript{μ}/m\textsubscript{τ} & (b₂ − D\textsubscript{bulk})/\textbar PSL(2,7)\textbar{} = 10/168 = \textbf{5/84} = 0.059524 & R8 & integer; approximate (R8); \textbf{proximity}; pull +20.8σ \\
\end{longtable}}

Four relations bear on two independent mass ratios, and they are mutually incompatible: with 27\textsuperscript{φ} and 3477, the Koide parameter is 0.666546, not 2/3. None is selected (main §6.1). Earlier versions added that 3477 = 3 × 19 × 61 ``factorizes into framework constants''; 19 is not a framework constant, and a factorization is not a derivation. They also cited a ``period-integral cross-confirmation'' of 3477 from associative 3-cycle volumes, ln(m\textsubscript{τ}/m\textsubscript{e}) = 8.154; that number is the logarithm of the measured ratio, used as the target of the volume assignment (S3, entries C17--C19), so the agreement is the target returning, not an independent confirmation. Both statements are withdrawn.

\begin{center}\rule{\linewidth}{0.5pt}\end{center}

\section{S2.4 Quarks, CKM and boson masses}\label{s2.4-quarks-ckm-and-boson-masses}

{\footnotesize\setlength{\tabcolsep}{3pt}
\begin{longtable}[]{@{}
  >{\raggedright\arraybackslash}p{(\linewidth - 8\tabcolsep) * \real{0.0870}}
  >{\raggedright\arraybackslash}p{(\linewidth - 8\tabcolsep) * \real{0.2174}}
  >{\raggedright\arraybackslash}p{(\linewidth - 8\tabcolsep) * \real{0.2391}}
  >{\raggedright\arraybackslash}p{(\linewidth - 8\tabcolsep) * \real{0.3913}}
  >{\raggedright\arraybackslash}p{(\linewidth - 8\tabcolsep) * \real{0.0652}}@{}}
\toprule\noalign{}
\begin{minipage}[b]{\linewidth}\raggedright
ID
\end{minipage} & \begin{minipage}[b]{\linewidth}\raggedright
Relation
\end{minipage} & \begin{minipage}[b]{\linewidth}\raggedright
Arithmetic
\end{minipage} & \begin{minipage}[b]{\linewidth}\raggedright
Choices declared
\end{minipage} & \begin{minipage}[b]{\linewidth}\raggedright
Status (main §4.1)
\end{minipage} \\
\midrule\noalign{}
\endhead
\bottomrule\noalign{}
\endlastfoot
A08 & m\textsubscript{s}/m\textsubscript{d} & p₂² × Weyl = 4 × 5 = \textbf{20} & the PDG gives a range (17--22), not a measurement & integer; inside the PDG range \\
A09 & m\textsubscript{c}/m\textsubscript{s} & (dim E₈ − p₂)/b₂ = 246/21 = \textbf{82/7} = 11.7143 & n/a & integer; pull −1.6σ \\
A10 & m\textsubscript{b}/m\textsubscript{t} & 1/(2b₂) = \textbf{1/42} = 0.023810 & \textbf{mass scheme} (R7): m\textsubscript{b} in MS-bar at m\textsubscript{b} over the directly measured m\textsubscript{t} & integer; \textbf{proximity} (in the scheme declared in R7); pull −10.2σ \\
A11 & m\textsubscript{u}/m\textsubscript{d} & (1 + dim E₆)/\textbar PSL(2,7)\textbar{} = \textbf{79/168} = 0.470238 & additive 1; the PDG quotes the uncertainty at 90\% CL, divided by 1.645 to 1σ (S3) & integer; pull +0.7σ \\
A16 & sin θ₁₂ (CKM) = λ & fund(E₇)/dim E₈ = 56/248 = \textbf{7/31} = 0.225806 & compared with sin θ₁₂, not sin²θ₁₂; R12 & integer; approximate (R12); pull +0.9σ \\
A17 & A (Wolfenstein) & (Weyl + dim E₆)/H* = \textbf{83/99} = 0.838384 & R12 & integer; approximate (R12); pull +0.7σ \\
A18 & sin θ₂₃ (CKM) = Aλ² & dim K₇/\textbar PSL(2,7)\textbar{} = 7/168 = \textbf{1/24} = 0.041667 & compared with sin θ₂₃, not sin²θ₂₃; R12 & integer; approximate (R12); pull −0.3σ \\
A19 & m\textsubscript{H}/m\textsubscript{t} & fund(E₇)/b₃ = 56/77 = \textbf{8/11} = 0.727273 & m\textsubscript{t} from direct measurements & integer; pull +1.8σ \\
A20 & m\textsubscript{H}/m\textsubscript{W} & (N\textsubscript{gen} + dim E₆)/dim F₄ = \textbf{81/52} = 1.557692 & n/a & integer; pull +0.5σ \\
A21 & m\textsubscript{W}/m\textsubscript{Z} & (2b₂ − Weyl)/(2b₂) = \textbf{37/42} = 0.880952 & n/a & integer; \textbf{tension}; pull −2.7σ \\
\end{longtable}}

\textbf{The CKM sector is overdetermined (R12).} The three relations bear on two Wolfenstein parameters, since sin θ₂₃ = Aλ². With λ = 7/31 and A = 83/99, Aλ² = 4067/95139 = 0.042748, while A18 gives 1/24 = 0.041667: the three relations are 2.6\% apart as exact equalities, and at most two can hold. None is selected (main §4.2).

\textbf{Labels corrected.} Earlier versions labelled A16 and A18 ``sin²θ₁₂(CKM)'' and ``sin²θ₂₃(CKM)'' while comparing them with λ = sin θ₁₂ and Aλ² = sin θ₂₃. The values are sines; the labels are corrected (S3.3).

\textbf{On the constant 42.} Earlier versions presented the appearance of 42 = 2b₂ in both m\textsubscript{b}/m\textsubscript{t} = 1/42 and Ω\textsubscript{DM}/Ω\textsubscript{b} = (1 + 42)/8 as ``connecting quark physics to large-scale structure''. Two formulas that share an integer share a choice of vocabulary; no connection between the quantities is implied, and the sentence is withdrawn.

\textbf{Dependence noted, not counted.} sin²θ\textsubscript{W} in MS-bar (A01) and m\textsubscript{W}/m\textsubscript{Z} (A21) are related in the Standard Model through radiative corrections; the two relations are not independent tests in the strict sense.

\begin{center}\rule{\linewidth}{0.5pt}\end{center}

\section{S2.5 Neutrino mixing}\label{s2.5-neutrino-mixing}

{\footnotesize\setlength{\tabcolsep}{3pt}
\begin{longtable}[]{@{}
  >{\raggedright\arraybackslash}p{(\linewidth - 8\tabcolsep) * \real{0.0870}}
  >{\raggedright\arraybackslash}p{(\linewidth - 8\tabcolsep) * \real{0.2174}}
  >{\raggedright\arraybackslash}p{(\linewidth - 8\tabcolsep) * \real{0.2391}}
  >{\raggedright\arraybackslash}p{(\linewidth - 8\tabcolsep) * \real{0.3913}}
  >{\raggedright\arraybackslash}p{(\linewidth - 8\tabcolsep) * \real{0.0652}}@{}}
\toprule\noalign{}
\begin{minipage}[b]{\linewidth}\raggedright
ID
\end{minipage} & \begin{minipage}[b]{\linewidth}\raggedright
Relation
\end{minipage} & \begin{minipage}[b]{\linewidth}\raggedright
Arithmetic
\end{minipage} & \begin{minipage}[b]{\linewidth}\raggedright
Choices declared
\end{minipage} & \begin{minipage}[b]{\linewidth}\raggedright
Status (main §4.1)
\end{minipage} \\
\midrule\noalign{}
\endhead
\bottomrule\noalign{}
\endlastfoot
A12 & sin²θ₁₂ (PMNS) & (1 + N\textsubscript{gen})/α\textsubscript{sum} = \textbf{4/13} = 0.307692 & additive 1 & integer; pull −0.2σ \\
A13 & sin²θ₂₃ (PMNS) & (D\textsubscript{bulk} − Weyl)/D\textsubscript{bulk} = \textbf{6/11} = 0.545455 & octant (main §4.3) & integer; \textbf{tension} (octant; inside the NuFIT 3σ range, main §4.3); pull +4.4σ, not Gaussian (main §4.3) \\
A14 & sin²θ₁₃ (PMNS) & D\textsubscript{bulk}/dim(E₈×E₈) = \textbf{11/496} = 0.022177 & n/a & integer; pull −0.5σ \\
A15 & δ\textsubscript{CP} & dim K₇ · dim G₂ + H* = 98 + 99 = \textbf{197}, read as the phase 2π · 197/360 & \textbf{unit} (R6) & \textbf{unit assignment} (R6); pull −0.4σ \\
\end{longtable}}

\textbf{The δ\textsubscript{CP} rule.} The prediction is the phase 2π · 197/360, a point value; the unit is a dictionary assignment (R6), and rule D1 of main §11 applies: a measurement whose 3σ interval excludes 197° rejects the relation and the unit assignment together, with no repair clause and no theoretical tolerance added afterwards. The compactification factor 62/69 of earlier versions is addressed in S2.8.

\begin{center}\rule{\linewidth}{0.5pt}\end{center}

\section{S2.6 Cosmology (overdetermined: R11)}\label{s2.6-cosmology-overdetermined-r11}

{\footnotesize\setlength{\tabcolsep}{3pt}
\begin{longtable}[]{@{}
  >{\raggedright\arraybackslash}p{(\linewidth - 8\tabcolsep) * \real{0.0870}}
  >{\raggedright\arraybackslash}p{(\linewidth - 8\tabcolsep) * \real{0.2174}}
  >{\raggedright\arraybackslash}p{(\linewidth - 8\tabcolsep) * \real{0.2391}}
  >{\raggedright\arraybackslash}p{(\linewidth - 8\tabcolsep) * \real{0.3913}}
  >{\raggedright\arraybackslash}p{(\linewidth - 8\tabcolsep) * \real{0.0652}}@{}}
\toprule\noalign{}
\begin{minipage}[b]{\linewidth}\raggedright
ID
\end{minipage} & \begin{minipage}[b]{\linewidth}\raggedright
Relation
\end{minipage} & \begin{minipage}[b]{\linewidth}\raggedright
Arithmetic
\end{minipage} & \begin{minipage}[b]{\linewidth}\raggedright
Choices declared
\end{minipage} & \begin{minipage}[b]{\linewidth}\raggedright
Status (main §4.1)
\end{minipage} \\
\midrule\noalign{}
\endhead
\bottomrule\noalign{}
\endlastfoot
A23 & Ω\textsubscript{DE} = Ω\textsubscript{Λ} & ln 2 · (b₂ + b₃)/H* = ln 2 · 98/99 = \textbf{0.686146} & operation ln 2; R11 & ln 2; approximate (R11); pull +0.2σ \\
A24 & n\textsubscript{s} & ζ(D\textsubscript{bulk})/ζ(Weyl) = \textbf{ζ(11)/ζ(5)} = 0.964864 & operation ζ & ζ values; pull 0.0σ \\
A25 & Ω\textsubscript{DM}/Ω\textsubscript{b} & (1 + 2b₂)/rank E₈ = \textbf{43/8} = 5.375 & additive 1; R11 & integer; approximate (R11); pull +0.1σ \\
A26 & Ω\textsubscript{c}/Ω\textsubscript{Λ} & det\textsubscript{num}/\textbar PSL(2,7)\textbar{} = \textbf{65/168} = 0.386905 & R2; R11 & integer; det(g); approximate (R11); pull 0.0σ \\
A27 & Ω\textsubscript{Λ}/Ω\textsubscript{m} & (dim G₂ + H*)/dim F₄ = \textbf{113/52} = 2.173077 & R11 & integer; approximate (R11); pull 0.0σ \\
A28 & h & (\textbar PSL(2,7)\textbar{} − 1)/dim E₈ = \textbf{167/248} = 0.673387 & additive 1 & integer; pull 0.0σ \\
A29 & Ω\textsubscript{b}/Ω\textsubscript{m} & Weyl/det\textsubscript{den} = \textbf{5/32} = 0.15625 & R2; R11 & integer; det(g); approximate (R11); pull −0.1σ \\
A30 & Ω\textsubscript{c}/Ω\textsubscript{m} & (det\textsubscript{den} − Weyl)/det\textsubscript{den} = \textbf{27/32} = 0.84375 & R2; R11 & integer; det(g); approximate (R11); \textbf{tension}; pull +2.5σ \\
A31 & σ₈ & (p₂ + det\textsubscript{den})/(2b₂) = 34/42 = \textbf{17/21} = 0.809524 & R2 & integer; det(g); pull −0.3σ \\
A32 & Y\textsubscript{p} & (1 + dim G₂)/κ\textsubscript{T}⁻¹ = \textbf{15/61} = 0.245902 & additive 1; R3 & integer; κ\textsubscript{T}; pull +0.3σ \\
A33 & Ω\textsubscript{m}/Ω\textsubscript{Λ} & (1 + dim J₃(\ensuremath{\mathbb{O}}))/κ\textsubscript{T}⁻¹ = \textbf{28/61} = 0.459016 & additive 1; R3; R11 & integer; κ\textsubscript{T}; approximate (R11); pull −0.1σ \\
\end{longtable}}

\textbf{The seven density relations bear on two independent densities} (flat ΛCDM: Ω\textsubscript{b} + Ω\textsubscript{c} + Ω\textsubscript{Λ} = 1) and are mutually incompatible as exact equalities (main §4.2):

\begin{longtable}[]{@{}
  >{\raggedright\arraybackslash}p{(\linewidth - 2\tabcolsep) * \real{0.2381}}
  >{\raggedright\arraybackslash}p{(\linewidth - 2\tabcolsep) * \real{0.7619}}@{}}
\toprule\noalign{}
\begin{minipage}[b]{\linewidth}\raggedright
Quantity
\end{minipage} & \begin{minipage}[b]{\linewidth}\raggedright
Values implied by the relations
\end{minipage} \\
\midrule\noalign{}
\endhead
\bottomrule\noalign{}
\endlastfoot
Ω\textsubscript{Λ}/Ω\textsubscript{m} & 113/52 = 2.1731 (A27); 61/28 = 2.1786 (inverse of A33); 2.1862 (A23 with flatness) \\
Ω\textsubscript{c}/Ω\textsubscript{b} & 43/8 = 5.375 (A25); 27/5 = 5.400 (A30/A29) \\
Ω\textsubscript{c}/Ω\textsubscript{Λ} & 65/168 = 0.3869 (A26); 0.3873 (A30 × A33); 0.3883 (A30/A27) \\
\end{longtable}

A27 and A33 are the same quantity, written with two formulas. The disagreements are below the Planck uncertainties of these ratios, which is why six of the seven pass at 2σ; Ω\textsubscript{c}/Ω\textsubscript{m}, whose uncertainty is much smaller (0.0019), is in tension at +2.5σ. The seven cannot hold together; none is selected (R11). Some are redundant and some subsets are compatible: A29 and A30 are equivalent when neutrinos are neglected (Ω\textsubscript{m} = Ω\textsubscript{b} + Ω\textsubscript{c}, so that 5/32 + 27/32 = 1), with or without flatness, and the largest compatible subsets have three relations, A29 and A30 with one of A23, A26, A27 or A33 (checked by solving each subset for (Ω\textsubscript{b}, Ω\textsubscript{c})). Earlier drafts of this version said that at most two of the seven can be exact; A23, A29 and A30 hold together, and the statement is withdrawn.

\textbf{Withdrawn.} Earlier versions also gave Ω\textsubscript{m} = Ω\textsubscript{DE}/√Weyl = 0.3068 (2.7\% from Planck) as a ``derived'' relation, and described ln 2 as a ``binary information factor'' of the moduli space. The first is an eighth density relation, incompatible with the others; the second names an operation without a computation behind it. Both are withdrawn. The statement that Ω\textsubscript{DM}/Ω\textsubscript{b} = 43/8 is ``exact'' (from a dataset value 5.375 ± 0.1 that matched no Planck column) is withdrawn; against the Planck 2018 chains it is 5.364 ± 0.077.

\begin{center}\rule{\linewidth}{0.5pt}\end{center}

\section{S2.7 Vocabulary freedom: alternative formulas in the corpus}\label{s2.7-vocabulary-freedom-alternative-formulas-in-the-corpus}

Earlier versions counted, for each relation, the number of distinct expressions in the constants that reduce to the same value (for example 14 for sin²θ\textsubscript{W} = 3/13, 20 for Q\textsubscript{Koide} = 2/3, 21 for m\textsubscript{b}/m\textsubscript{t} = 1/42), and read a large count as ``structural robustness'': a value reached by many paths was said to be less likely a coincidence. Version 3.6 reverses that reading (main §7). A value reached by many paths through twenty constants is a value the vocabulary reaches easily; the count measures the vocabulary's freedom (R4, R5). It is not what the look-elsewhere test of main §7 prices: the published values 3/13 and 1/42 lie outside their 2σ windows, and the 30 targets that every competitor reaches are already reached by the grammar's universal leaves, with no constant of K₇ at all (S4.2.2). The counts are kept as such a measure:

\begin{longtable}[]{@{}
  >{\raggedright\arraybackslash}p{(\linewidth - 4\tabcolsep) * \real{0.1961}}
  >{\raggedright\arraybackslash}p{(\linewidth - 4\tabcolsep) * \real{0.6078}}
  >{\raggedright\arraybackslash}p{(\linewidth - 4\tabcolsep) * \real{0.1961}}@{}}
\toprule\noalign{}
\begin{minipage}[b]{\linewidth}\raggedright
Relation
\end{minipage} & \begin{minipage}[b]{\linewidth}\raggedright
Distinct expressions recorded
\end{minipage} & \begin{minipage}[b]{\linewidth}\raggedright
Examples
\end{minipage} \\
\midrule\noalign{}
\endhead
\bottomrule\noalign{}
\endlastfoot
sin²θ\textsubscript{W} = 3/13 & 14 & N\textsubscript{gen}/α\textsubscript{sum}; b₂/(b₃ + dim G₂); dim J₃(\ensuremath{\mathbb{O}})/(dim F₄ + det\textsubscript{num}); (1 + dim G₂)/det\textsubscript{num} \\
Q\textsubscript{Koide} = 2/3 & 20 & p₂/N\textsubscript{gen}; dim G₂/b₂; dim F₄/dim E₆; rank E₈/(Weyl + dim K₇) \\
m\textsubscript{b}/m\textsubscript{t} = 1/42 & 21 & 1/(2b₂); (1 + N\textsubscript{gen})/\textbar PSL(2,7)\textbar; p₂/(dim K₇ + b₃); N\textsubscript{gen}/(dim J₃(\ensuremath{\mathbb{O}}) + H*) \\
sin²θ₁₂ (PMNS) = 4/13 & 28 & n/a \\
m\textsubscript{H}/m\textsubscript{t} = 8/11 & 19 & n/a \\
m\textsubscript{u}/m\textsubscript{d} = 79/168, m\textsubscript{H}/m\textsubscript{W} = 81/52 & 1 each & n/a \\
\end{longtable}

\textbf{Second formulas for the same observable.} The corpus also contains formulas that give a \emph{different} value for the same observable, chosen at different times:

\begin{longtable}[]{@{}
  >{\raggedright\arraybackslash}p{(\linewidth - 4\tabcolsep) * \real{0.2069}}
  >{\raggedright\arraybackslash}p{(\linewidth - 4\tabcolsep) * \real{0.2931}}
  >{\raggedright\arraybackslash}p{(\linewidth - 4\tabcolsep) * \real{0.5000}}@{}}
\toprule\noalign{}
\begin{minipage}[b]{\linewidth}\raggedright
Observable
\end{minipage} & \begin{minipage}[b]{\linewidth}\raggedright
Dataset formula
\end{minipage} & \begin{minipage}[b]{\linewidth}\raggedright
Other formula in the corpus
\end{minipage} \\
\midrule\noalign{}
\endhead
\bottomrule\noalign{}
\endlastfoot
λ\textsubscript{H} & √17/32 = 0.12885 & b₂(M₁)/(b₃ + b₂(M₂)) = 11/87 = 0.12644 (TCS building blocks) \\
Ω\textsubscript{DE} & ln 2 · 98/99 = 0.68615 & 1 − χ(K3)/b₃ = 53/77 = 0.68831 \\
m\textsubscript{c}/m\textsubscript{s} & 246/21 = 11.714 & b₂(M₁) + dim K₇/b₂(M₂) = 117/10 = 11.700 \\
Ω\textsubscript{Λ}/Ω\textsubscript{m} & 113/52 = 2.1731 & (det\textsubscript{den} − dim K₇)/D\textsubscript{bulk} = 25/11 = 2.2727; 61/28 (A33 inverted) \\
θ₁₃ & sin²θ₁₃ = 11/496 & θ₁₃ = π/b₂ = 8.571°; θ₁₃ = 94/11 = 8.545° read in degrees \\
θ₂₃ & sin²θ₂₃ = 6/11 (θ = 47.6°) & θ₂₃ = arcsin(75/99) = 49.25° (sin²θ₂₃ = 0.574) \\
θ₁₂ & sin²θ₁₂ = 4/13 & θ₁₂ = arctan √((2π/25)/(511/884)) = 33.40° \\
δ\textsubscript{CP} & 197 & 197 × 62/69 = 177.01 (S2.8) \\
\end{longtable}

The ``TCS building-block'' formulas use b₂(M₁) = 11 and b₂(M₂) = 10, Betti numbers of building blocks whose identification is open (main §9). Several alternatives were introduced after a release of the data had moved: the Ω\textsubscript{DE} alternative was preferred against one Planck release and not the other, and the θ₁₃ alternative ``in degrees'' is a second unit assignment of the kind R6 registers. None of these is part of the 33; they are listed because they show, on the framework's own record, how many formulas the vocabulary offers for one observable and how the choice among them was made.

\textbf{Patterns.} The Fibonacci and Lucas numbers among the constants (2, 3, 5, 8, 13, 21; 7, 11) and the identities among the integers (dim G₂ = p₂ · dim K₇, b₂ = N\textsubscript{gen} · dim K₇, \textbar PSL(2,7)\textbar{} = rank E₈ · b₂, b₃ + dim G₂ = dim K₇ · α\textsubscript{sum}) are arithmetic facts about the chosen integers. Earlier versions called them an ``algebraic web'' and a ``Fibonacci--Lucas embedding'' reflecting icosahedral geometry. Small integers satisfy many such identities; no physical consequence follows from these, and they are not evidence.

\begin{center}\rule{\linewidth}{0.5pt}\end{center}

\section{S2.8 Withdrawn readings}\label{s2.8-withdrawn-readings}

\begin{longtable}[]{@{}
  >{\raggedright\arraybackslash}p{(\linewidth - 2\tabcolsep) * \real{0.4872}}
  >{\raggedright\arraybackslash}p{(\linewidth - 2\tabcolsep) * \real{0.5128}}@{}}
\toprule\noalign{}
\begin{minipage}[b]{\linewidth}\raggedright
Earlier statement
\end{minipage} & \begin{minipage}[b]{\linewidth}\raggedright
Why it is withdrawn
\end{minipage} \\
\midrule\noalign{}
\endhead
\bottomrule\noalign{}
\endlastfoot
``Complete derivations'', ``proof'', ``VERIFIED'', ``TOPOLOGICAL'' as statuses of the relations & the arithmetic is exact; the assignment is not derived (S2.0) \\
N\textsubscript{gen} = 3 by the Atiyah--Singer index theorem & no operator is specified; the expression is the integer identity of S2.1 \\
``All 33 predictions use only topological invariants; no fitting'' & eight use the imposed normalization or κ\textsubscript{T} (R2, R3); the formulas were chosen with the values known (R4) \\
Mean deviation 0.73\%, ``11 exact matches'', deviation histograms & computed against an earlier dataset; a mean deviation measures closeness, which is not distinctive (main §4.2, §7) \\
Statistical uniqueness: 192 349 alternatives, none better, P \textless{} 5 × 10⁻⁶, gauge-group and holonomy rankings & the alternatives were scored with the framework's own formulas, not allowed to search; replaced by the preregistered null of main §7 and S4 \\
``Literature uniqueness'' of (21, 77) & absence from a catalogue is not selection; the selection audit is main §10 \\
Results ``from analysis of the NK-certified metric'': V\textsubscript{min} formula, U(1)² symmetry, origin of the SM gauge group from triality & properties of an approximate seam metric (main §9); not used by any relation, and the gauge-group statement is not a derivation (main §2) \\
κ\textsubscript{T} = 1/61 as a ``topological torsion capacity''; 61 = dim F₄ + N\textsubscript{gen}² & an integer combination without a physical name (R3); a second expression for 61 is another instance of S2.7 \\
det(g) = 65/32 ``consistent with the topological formula'' via three algebraic paths & imposed normalization (R2); the metric optimization was constrained to this target \\
3477 ``factorizes into framework constants''; associative-volume ``cross-confirmation'' & S2.3 \\
Ω\textsubscript{m} = Ω\textsubscript{DE}/√5, Ω\textsubscript{DM}/Ω\textsubscript{b} ``exact'', ln 2 as an information factor & S2.6 \\
\end{longtable}

\textbf{Appendix F of earlier versions: the compactification factor 62/69.} A residual search after the NuFIT 6.0 release found 197 × 62/69 = 177.01, matching that release's central value, and earlier versions kept it as a ``contingency plan'': if DUNE measured δ\textsubscript{CP} near 177°, the framework would adopt 177.01 as its prediction and read 197 as a ``bare'' value. That plan is a repair clause, and version 3.6 has none (main §11, rule D1). It is withdrawn without condition: if a measurement excludes 197° at 3σ, the relation δ\textsubscript{CP} = 197 and its unit assignment are rejected, and no factor found afterwards replaces them. The factor 62/69 = 248/276 is recorded in S2.7 as one more formula the vocabulary offers.

\begin{center}\rule{\linewidth}{0.5pt}\end{center}


\bigskip
% GÉNÉRÉ par gen_bib_3_6.py depuis physics_main.md : ne pas éditer.

\begin{thebibliography}{99}

\bibitem{1} Particle Data Group, F. Takahashi et al., Int. J. Mod. Phys. A 41, 2630011 (2026)

\bibitem{2} S. Weinberg, Phys. Rev.~D 13, 974 (1976)

\bibitem{3} Planck Collaboration, A\&A 641, A6 (2020); chains: Planck Legacy Archive, COM\_CosmoParams\_base-plikHM-TTTEEE-lowl-lowE\_R3.00

\bibitem{4} A.G. Riess et al., ApJL 934, L7 (2022)

\bibitem{5} C.D. Froggatt, H.B. Nielsen, Nucl. Phys. B 147, 277 (1979)

\bibitem{6} Y. Koide, Lett. Nuovo Cim. 34, 201 (1982)

\bibitem{7} C. Furey, \emph{Standard Model Physics from an Algebra?}, PhD thesis, U. Waterloo (2015)

\bibitem{8} C. Furey, M.J. Hughes, Phys. Lett. B 831, 137186 (2022)

\bibitem{9} R. Wilson, arXiv:2404.18938 (2024)

\bibitem{10} T.P. Singh et al., arXiv:2206.06911v3 (2024)

\bibitem{11} B.S. Acharya, S. Gukov, Phys. Rep.~392, 121 (2004)

\bibitem{12} L. Foscolo et al., Duke Math. J. 170, 3 (2021)

\bibitem{13} A. Kovalev, ``Coassociative K3 fibrations of compact G₂-manifolds,'' arXiv:math/0511150 (2005)

\bibitem{14} M. Haskins, J. Nordstrom, arXiv:1809.09083 (2022)

\bibitem{15} G. Bera, ``Associative submanifolds in twisted connected sum G₂-manifolds,'' arXiv:2209.00156 (2022)

\bibitem{16} D. Crowley, S. Goette, J. Nordstrom, Invent. Math. (2025)

\bibitem{17} T. Dray, C.A. Manogue, \emph{The Geometry of the Octonions}, World Scientific (2015)

\bibitem{18} J.F. Adams, \emph{Lectures on Exceptional Lie Groups}, U. Chicago Press (1996)

\bibitem{19} D.J. Gross et al., Nucl. Phys. B 256, 253 (1985)

\bibitem{20} D.D. Joyce, \emph{Compact Manifolds with Special Holonomy}, Oxford U. Press (2000)

\bibitem{21} A. Kovalev, J. Reine Angew. Math. 565, 125 (2003)

\bibitem{22} A. Corti et al., Duke Math. J. 164, 1971 (2015)

\bibitem{23} R. Harvey, H.B. Lawson, Acta Math. 148, 47 (1982)

\bibitem{24} B.S. Acharya, Class. Quant. Grav. 19, 5619 (2002)

\bibitem{25} B.S. Acharya, E. Witten, arXiv:hep-th/0109152 (2001)

\bibitem{26} The NOvA Collaboration, The T2K Collaboration, ``Joint neutrino oscillation analysis from the T2K and NOvA experiments,'' Nature 646, 818-824 (2025), doi:10.1038/s41586-025-09599-3, arXiv:2510.19888

\bibitem{27} I. Esteban, M.C. Gonzalez-Garcia, M. Maltoni, I. Martinez-Soler, J.P. Pinheiro, T. Schwetz, ``NuFit-6.0: updated global analysis of three-flavor neutrino oscillations,'' JHEP 12 (2024) 216, arXiv:2410.05380; NuFIT 6.1: three-neutrino fit based on data available in November 2025, www.nu-fit.org

\bibitem{28} L. de Moura, S. Ullrich, CADE 28, 625 (2021)

\bibitem{29} mathlib Community, github.com/leanprover-community/mathlib4

\bibitem{30} DUNE Collaboration, FERMILAB-TM-2696 (2020)

\bibitem{31} DUNE Collaboration, arXiv:2103.04797 (2021)

\bibitem{32} E. Witten, ``Strong coupling expansion of Calabi-Yau compactification,'' Nucl. Phys. B 471, 135 (1996)

\bibitem{33} B.S. Acharya et al., Phys. Rev.~D 76, 126010 (2007)

\bibitem{34} M. Atiyah, E. Witten, Adv. Theor. Math. Phys. 6, 1 (2002)

\bibitem{35} G. Kane, \emph{String Theory and the Real World} (2017)

\bibitem{36} J. Distler, S. Garibaldi, Commun. Math. Phys. 298, 419 (2010)

\bibitem{37} J.C. Baez, ``Octonions and the Standard Model,'' math.ucr.edu/home/baez/standard/ (2020-2025)

\bibitem{38} Springer Nature, ``Artificial intelligence (AI) policy,'' www.springernature.com/gp/policies (2024)

\bibitem{39} J.A. Wheeler, ``Information, physics, quantum: the search for links,'' in \emph{Complexity, Entropy, and the Physics of Information} (W.H. Zurek, ed.), Addison-Wesley (1990), pp.~3--28.

\bibitem{40} J. Worrall, ``Structural realism: the best of both worlds?,'' \emph{Dialectica} \textbf{43}, 99--124 (1989).

\bibitem{41} J. Ladyman, D. Ross, \emph{Every Thing Must Go: Metaphysics Naturalized}, Oxford University Press (2007).

\bibitem{42} R. Pinčák, A. Pigazzini, M. Pudlák, E. Bartoš, ``Geometric origin of a stable black hole remnant from torsion in G₂-manifold geometry,'' Gen.~Rel. Grav. \textbf{58}, 29 (2026), doi:10.1007/s10714-026-03528-z

\bibitem{43} D. Joyce, S. Karigiannis, ``A new construction of compact torsion-free G₂-manifolds by gluing families of Eguchi-Hanson spaces,'' \emph{J. Differential Geom.} (2021), arXiv:1707.09325 (2017).

\bibitem{44} S. Mukai, ``Finite groups of automorphisms of K3 surfaces and the Mathieu group,'' \emph{Invent. Math.} \textbf{94}, 183--221 (1988).

\bibitem{45} A. Garbagnati, A. Sarti, ``Symplectic automorphisms of prime order on K3 surfaces,'' \emph{J. Algebra} \textbf{318}, 323--350 (2007), arXiv:math/0603742.

\bibitem{46} LEP2 SUSY Working Group (ALEPH, DELPHI, L3, OPAL), ``Combined LEP chargino results, up to 208 GeV for large m0,'' note LEPSUSYWG/01-03.1 (2001), lepsusy.web.cern.ch.

\bibitem{47} Fermilab News, ``Fermilab marks major milestone for world-leading DUNE experiment,'' 7 May 2026, news.fnal.gov.

\bibitem{48} J. Chen, H. Hong, ``Intermediate curvature and splitting theorem,'' arXiv:2604.26529 (2026). (Cited in Supplement S1, §8.4.)

\bibitem{49} M. Tristram et al., ``Cosmological parameters derived from the final Planck data release (PR4),'' A\&A 682, A37 (2024). (Cited in Supplement S3; not used in the dataset.)

\bibitem{50} E. Tiesinga, P.J. Mohr, D.B. Newell, B.N. Taylor, ``CODATA recommended values of the fundamental physical constants: 2022,'' Rev.~Mod. Phys. 97, 025002 (2025). (Cited in Supplement S3.)

\bibitem{51} V. V. Nikulin, ``Integer symmetric bilinear forms and some of their geometric applications,'' Izv. Akad. Nauk SSSR Ser. Mat. 43 (1979); English transl. Math. USSR Izv. 14 (1980).

\bibitem{52} Belle II Collaboration, ``Measurement of the τ-lepton mass with the Belle II experiment,'' Phys. Rev.~D 108, 032006 (2023), doi:10.1103/PhysRevD.108.032006, arXiv:2305.19116.

\bibitem{53} S. K. Donaldson, ``Adiabatic limits of co-associative Kovalev--Lefschetz fibrations,'' in Algebra, Geometry, and Physics in the 21st Century, Progr. Math. 324, Birkhäuser (2017), pp.~1--29, arXiv:1603.08391.

\bibitem{54} C. Camere, ``Lattice polarized irreducible holomorphic symplectic manifolds,'' Ann. Inst. Fourier 66 (2016), 687--709, arXiv:1312.3224.

\bibitem{55} C. Zhou, Z. Zhou, ``Algebraic Stability and Cosmological Structure E: The Arithmetic Necessity of Three Generations of Fermions,'' Zenodo, doi:10.5281/zenodo.18626667 (2026).

\bibitem{56} Particle Data Group, S. Navas et al., ``Review of Particle Physics,'' Phys. Rev.~D 110, 030001 (2024), doi:10.1103/PhysRevD.110.030001 (cited in S3 only)

\bibitem{57} CDF Collaboration, T. Aaltonen et al., ``High-precision measurement of the W boson mass with the CDF II detector,'' Science 376, 170 (2022), doi:10.1126/science.abk1781

\bibitem{58} JUNO Collaboration, ``First measurement of reactor neutrino oscillations at JUNO,'' arXiv:2511.14593 (2025).

\bibitem{59} JUNO Collaboration, ``Sub-percent precision measurement of neutrino oscillation parameters with JUNO,'' Chin. Phys. C 46, 123001 (2022), arXiv:2204.13249.

\bibitem{60} D. Buttazzo, G. Degrassi, P.P. Giardino, G.F. Giudice, F. Sala, A. Salvio, A. Strumia, ``Investigating the near-criticality of the Higgs boson,'' JHEP 12 (2013) 089, arXiv:1307.3536.

\end{thebibliography}


\end{document}
